\chapter{Berkovich spectra and analytification}
\index{Berkovich space}\index{analytification!Berkovich}
\label{chap:berkovich}

The Zariski topology remembers specialization very well, but it is too coarse
for potential theory. The topology on the set of points of a non-Archimedean
field has the opposite defect: it has many clopen subsets and generally has
too few paths. Berkovich's construction adds seminorms, not merely points
over field extensions. The resulting spaces are Hausdorff, proper varieties
become compact, and degenerations acquire genuine polyhedral subsets.

This chapter develops the part of the construction needed in the rest of the
book. The reader only needs the valuation theory of \cref{chap:valuations}
and the language of models from \cref{chap:models}. We first study the
spectrum of a Banach algebra and prove its compactness in full,
then analytify schemes by gluing, establish the universal property of
analytification at the level we use it, explain centers on models, and
prove the complete classification of the points of the analytic affine
line over an algebraically closed field. Reduction maps are computed
explicitly on the projective line over \(\mathbf Q_p\) and over
\(\mathbf C(\!(t)\!)\). The final section emphasizes that the same
definition works over a trivially valued field, where we illustrate the
affine line and its rescaling rays. We construct the underlying
topological spaces carefully, but cite the general theory for structure
sheaves, relative properness, and analytic GAGA.

Throughout this chapter, \(K\) is a complete field with a non-Archimedean
absolute value \(|\cdot|\), possibly the trivial one. We use
\[
                         v(a)=-\log |a|.
\]
Thus \(v(0)=+\infty\). A multiplicative seminorm is allowed to have a
nonzero kernel.

\section{The spectrum of a Banach algebra}
\label{sec:berk-spectrum}

Let \(A\) be a commutative unital Banach \(K\)-algebra, with a
submultiplicative norm \(\|\cdot\|\).

\begin{definition}[Berkovich spectrum]
The \emph{Berkovich spectrum} \(\mathcal M(A)\) is the set of maps
\[
 |\cdot|_x:A\longrightarrow \mathbf R_{\geq 0}
\]
such that, for \(f,g\in A\),
\[
\begin{split}
 |0|_x&=0,\qquad |1|_x=1,\qquad |fg|_x=|f|_x|g|_x,\\
 |f+g|_x&\leq\max\{|f|_x,|g|_x\},\qquad
 |a|_x=|a|\quad(a\in K),\\
 |f|_x&\leq\|f\|.
\end{split}
\]
It is endowed with the weakest topology for which every evaluation map
\[
       x\longmapsto |f(x)|:=|f|_x,\qquad f\in A,
\]
is continuous.
\end{definition}

Some authors initially require only \( |f|_x\leq C\|f\| \) for one constant
\(C\). The following lemma shows that multiplicativity upgrades the weaker
inequality to the normalized one, and in fact to a bound by the spectral
seminorm.  Recall that the \emph{spectral seminorm}\index{spectral seminorm}
of \(f\in A\) is
\[
                     \rho(f)=\lim_{m\to\infty}\|f^m\|^{1/m}.
\]

\begin{lemma}\label{lem:berk-spectral-bound}
The limit defining \(\rho(f)\) exists and equals
\(\inf_{m\geq1}\|f^m\|^{1/m}\leq\|f\|\).  Moreover, if
\(|\cdot|_x:A\to\mathbf R_{\geq0}\) is multiplicative with \(|1|_x=1\) and
satisfies \(|f|_x\leq C\|f\|\) for all \(f\) and one constant \(C\geq1\),
then
\[
                     |f|_x\leq\rho(f)\leq\|f\|
                     \qquad(f\in A).
\]
\end{lemma}

\begin{proof}
Set \(a_m=\log\|f^m\|\in[-\infty,\infty)\), with the convention that all
assertions are trivial if some power of \(f\) is zero.  Submultiplicativity
of the norm gives \(a_{m+n}\leq a_m+a_n\).  Fekete's subadditivity argument
applies: fix \(n\geq1\), write \(m=qn+s\) with \(0\leq s<n\), and estimate
\[
 \frac{a_m}{m}\leq\frac{q\,a_n+a_s}{m}
 \xrightarrow[m\to\infty]{}\frac{a_n}{n},
\]
where \(a_0:=0\).  Hence \(\limsup_m a_m/m\leq\inf_n a_n/n\leq\liminf_m
a_m/m\), so the limit exists and equals the infimum; taking \(n=1\) gives
\(\rho(f)\leq\|f\|\).

For the second assertion, apply the assumed inequality to \(f^m\):
\[
 |f|_x^m=|f^m|_x\leq C\|f^m\|,
 \qquad\text{so}\qquad
 |f|_x\leq C^{1/m}\|f^m\|^{1/m}.
\]
Letting \(m\to\infty\) gives \(|f|_x\leq\rho(f)\).
\end{proof}

In particular, replacing the requirement \(|f|_x\le\|f\|\) in the
definition by boundedness with any constant changes nothing, and every
point of \(\mathcal M(A)\) automatically satisfies the finer inequality
\(|f|_x\le\rho(f)\).  Berkovich's spectral radius formula asserts that
\(\rho(f)=\max_{x\in\mathcal M(A)}|f(x)|\); we use only the elementary
inequality just proved and cite the formula, when needed, from
\cite[Theorem~1.3.1]{Berk12}.

\begin{proposition}\label{prop:berk-compact-spectrum}
The space \(\mathcal M(A)\) is compact and Hausdorff. If \(A\) is a nonzero
\(K\)-affinoid algebra, then \(\mathcal M(A)\) is nonempty.
\end{proposition}

\begin{proof}
We give the weak-\(*\) compactness argument in full.  Consider the product
\[
                     P=\prod_{f\in A}[0,\|f\|],
\]
equipped with the product topology; it is compact by Tychonoff's theorem
\cite[Chapter~5]{Kel75} and its topology is that of pointwise convergence
of \([0,\infty)\)-valued functions on \(A\) obeying the norm bound.  A
point of \(\mathcal M(A)\) is by definition such a function, so there is a
natural injection
\[
 \iota:\mathcal M(A)\longrightarrow P,
 \qquad
 \iota(x)=\bigl(|f|_x\bigr)_{f\in A}.
\]
Injectivity holds because a seminorm is its family of values.  The
topology of \(\mathcal M(A)\) was defined as the weakest one making every
evaluation \(x\mapsto|f|_x\) continuous; the subspace topology induced by
\(P\) has exactly the same description, since the projections
\(\mathrm{pr}_f\) generate the product topology.  Hence \(\iota\) is a
homeomorphism onto its image, and it suffices to prove that this image is
closed in \(P\).

Let \(\varphi\in P\) be adherent to \(\iota(\mathcal M(A))\); we check
each axiom.  Fix \(f,g\in A\) and \(a\in K\).  The maps
\[
 \varphi\longmapsto\varphi(fg)-\varphi(f)\varphi(g),
 \qquad
 \varphi\longmapsto\varphi(f+g)-\max\{\varphi(f),\varphi(g)\}
\]
are continuous on \(P\), because each involves only the three coordinates
indexed by \(f\), \(g\), and their sum or product, together with
 continuous operations on real numbers.  The conditions
 \[
 \begin{gathered}
 \varphi(0)=0,\qquad \varphi(1)=1,
 \qquad \varphi(fg)=\varphi(f)\varphi(g),\\
 \varphi(f+g)\leq\max\{\varphi(f),\varphi(g)\},
 \qquad \varphi(a\cdot1)=|a|
 \end{gathered}
 \]
are therefore equalities and weak inequalities between continuous
functions, and each cuts out a closed subset of \(P\).  The image of
\(\iota\) is the intersection of these closed subsets over all \(f\),
\(g\), and \(a\), hence is closed.  A closed subspace of a compact space
is compact, so \(\mathcal M(A)\) is compact.

For the Hausdorff property, let \(x\neq y\) in \(\mathcal M(A)\).  Some
\(f\in A\) has \(|f|_x\neq|f|_y\); choose disjoint open intervals
\(I\ni|f|_x\) and \(J\ni|f|_y\) in \(\mathbf R\).  The sets
\(\{z:|f|_z\in I\}\) and \(\{z:|f|_z\in J\}\) are disjoint open
neighborhoods.

Nonemptiness holds in much greater generality: the spectrum of every
nonzero commutative Banach ring is nonempty.  This is the non-Archimedean
spectral existence theorem, \cite[Theorem~1.2.1]{Berk12}; its proof
produces a bounded multiplicative seminorm through a maximal ideal of
\(A\) and an extension of the ground field, and we do not reproduce it.
A nonzero \(K\)-affinoid algebra is in particular a nonzero Banach ring,
which gives the stated case.
\end{proof}

Each point has both an algebraic residue field and an analytic residue field.
They should not be confused.

\begin{definition}[Support and completed residue field]
For \(x\in\mathcal M(A)\), set
\[
 \mathfrak p_x=\{f\in A:|f(x)|=0\}.
\]
This is a prime ideal, called the \emph{support} of \(x\). The seminorm
induces an absolute value on the fraction field of \(A/\mathfrak p_x\). Its
completion is denoted by \(\mathcal H(x)\) and called the \emph{completed
residue field} of \(x\).
\end{definition}

The point \(x\) can equivalently be viewed as the bounded homomorphism
\(A\to\mathcal H(x)\) together with the absolute value of
\(\mathcal H(x)\). This formulation makes residue fields and base change
natural.

\begin{example}\label{ex:berk-tate-disc}
The spectrum \(\mathcal M(K)\) consists of one point: a multiplicative
seminorm on a field extending \(|\cdot|\) has zero kernel and is determined
on \(K\) by the extension requirement. More interestingly,
for \(r>0\) let
\[
 K\langle r^{-1}T\rangle
 =\left\{\sum_{i\geq0}a_iT^i:\ |a_i|r^i\longrightarrow0\right\},
 \qquad
 \left\|\sum_i a_iT^i\right\|_r=\max_i|a_i|r^i.
\]
Its spectrum is the closed analytic disc of radius \(r\) about \(0\); this
identification is proved in \Cref{lem:berk-polydisc-exhaustion} below. The
norm itself defines a point: the seminorm
\[
             \left|\sum_i a_iT^i\right|_{\zeta_{0,r}}
             =\max_i |a_i|r^i
\]
is the \emph{Gauss point}\index{Gauss point} of that disc.  All axioms
except multiplicativity are immediate, and multiplicativity is the Gauss
lemma.  Here is the complete argument in the graded language of
\cref{chap:valuations}.  Write \(t=-\log r\) and
\[
 v_r\Bigl(\sum_i a_iT^i\Bigr)
 =-\log\Bigl\|\sum_i a_iT^i\Bigr\|_r
 =\min_i\{v(a_i)+it\}.
\]
The minimum is attained at finitely many indices, because
\(v(a_i)+it\to+\infty\).  Declare \(a_iT^i\) to have weight
\(v(a_i)+it\), and let \(\operatorname{in}(f)\) be the sum of the terms of
\(f\) of minimal weight \(v_r(f)\), viewed as an element of degree
\(v_r(f)\) in \(\operatorname{gr}_vK[\tau]\), where \(\tau\) is an
indeterminate of degree \(t\) representing the initial form of \(T\) and
each coefficient is replaced by its initial form in
\(\operatorname{gr}_vK\).  Collecting terms of least total weight in the
Cauchy product shows that either
\(\operatorname{in}(fg)=\operatorname{in}(f)\operatorname{in}(g)\), or the
polynomial product \(\operatorname{in}(f)\operatorname{in}(g)\) vanishes
and \(v_r(fg)>v_r(f)+v_r(g)\).  By
\cref{chap:valuations}, \(\operatorname{gr}_vK\) is a graded domain, so
the polynomial ring \(\operatorname{gr}_vK[\tau]\) is a domain and the
product of the two nonzero initial forms cannot vanish.  Hence
\(v_r(fg)=v_r(f)+v_r(g)\), which is multiplicativity of
\(\|\cdot\|_r\).
\end{example}

\section{Affinoid domains and gluing}
\label{sec:berk-affinoid-gluing}

The compact space \(\mathcal M(A)\) carries a sheaf of analytic functions.
One cannot obtain that sheaf by treating every ordinary open subset as an
affinoid chart. Instead one uses analytic domains and a Grothendieck
topology. We record one class of domains that is sufficient for intuition.

Suppose \(A\) is affinoid and \(f_1,\ldots,f_m,g\in A\) generate the unit
ideal. The associated \emph{rational domain} is
\[
 U=\{x\in\mathcal M(A): |f_i(x)|\leq |g(x)|,\ 1\leq i\leq m\}.
\]
It is represented by the affinoid algebra
\[
 A_U=A\langle T_1,\ldots,T_m\rangle/(gT_1-f_1,\ldots,gT_m-f_m),
\]
with the quotient Banach structure. Here
\(A\langle T_1,\ldots,T_m\rangle\) is the relative Tate
algebra\index{Tate algebra} of series \(\sum_\nu a_\nu T^\nu\) with
\(a_\nu\in A\) and \(\|a_\nu\|\to0\), normed by
\(\|\sum_\nu a_\nu T^\nu\|=\max_\nu\|a_\nu\|\).  The representation means
that bounded maps out of \(A\) for which the displayed inequalities hold
factor uniquely through this algebra. Rational domains form a basis for the
admissible topology, and Tate acyclicity makes their rings of functions
glue; for the sheaf-theoretic assertions we cite
\cite[\S2.2]{Berk12} and \cite{BoschLutkebohmert1993}.  At the level of
points, however, the meaning of ``represented'' admits a complete
elementary proof, which we now give because the same computation recurs
for polydisc exhaustions and for fibres of base change.

\begin{proposition}
\label{prop:berk-rational-domain}
Let \(A\) be a \(K\)-affinoid algebra, let \(f_1,\ldots,f_m,g\in A\)
generate the unit ideal, and let \(U\) and \(A_U\) be as above. The
bounded homomorphism \(A\to A_U\) induces a homeomorphism
\[
             \mathcal M(A_U)\xrightarrow{\ \sim\ }U.
\]
In particular \(U\) is a compact subset of \(\mathcal M(A)\), closed in
\(\mathcal M(A)\).
\end{proposition}

\begin{proof}
Write \(\varphi:A\to A_U\) for the structure map and
\(\pi:\mathcal M(A_U)\to\mathcal M(A)\), \(\pi(x)=x\circ\varphi\), for
the induced map; it is continuous because each evaluation
\(x\mapsto|\varphi(f)(x)|\) is.

\emph{The image lies in \(U\).}
Let \(x\in\mathcal M(A_U)\).  In \(A_U\) one has \(gT_i=f_i\), and
\(\|T_i\|\leq1\) in the Tate algebra, hence also in the quotient norm; by
\Cref{lem:berk-spectral-bound}, \(|T_i(x)|\leq1\).  Therefore
\[
 |f_i(x)|=|g(x)|\,|T_i(x)|\leq|g(x)|,
\]
so \(\pi(x)\in U\), provided we also know \(U\)-membership makes sense,
i.e.\ nothing further vanishes.  Note in passing that \(|g(x)|\neq0\):
choose \(a,b_1,\ldots,b_m\in A\) with \(ag+\sum_ib_if_i=1\); if
\(|g(x)|=0\), then every \(|f_i(x)|\leq|g(x)|=0\), and the ultrametric
inequality would give \(1=|1(x)|\leq\max\{|ag(x)|,|b_if_i(x)|\}=0\), a
contradiction.  The same argument shows \(|g(y)|\neq0\) for every
\(y\in U\).

\emph{Surjectivity.}
Let \(y\in U\), with character \(\chi_y:A\to\mathcal H(y)\).  Since
\(g(y)\neq0\) in the field \(\mathcal H(y)\), the elements
\(c_i=f_i(y)/g(y)\in\mathcal H(y)\) are defined, and \(|c_i|\leq1\) by the
definition of \(U\).  Define
\[
 \Phi:A\langle T_1,\ldots,T_m\rangle\longrightarrow\mathcal H(y),
 \qquad
 \Phi\Bigl(\sum_\nu a_\nu T^\nu\Bigr)=\sum_\nu a_\nu(y)\,c^\nu.
\]
The series on the right converges in the complete field
\(\mathcal H(y)\) because \(|a_\nu(y)c^\nu|\leq\|a_\nu\|\to0\), and
\(|\Phi(h)|\leq\|h\|\), so \(\Phi\) is a bounded homomorphism.  It kills
each \(gT_i-f_i\), hence factors through \(A_U\).  Composing with the
absolute value of \(\mathcal H(y)\) yields a point
\(x\in\mathcal M(A_U)\) with \(\pi(x)=y\).

\emph{Injectivity.}
Let \(x\in\mathcal M(A_U)\) with \(\pi(x)=y\), and let
\(\chi_x:A_U\to\mathcal H(x)\) be its character.  On the image of \(A\),
the seminorm of \(x\) equals that of \(y\); in particular
\(\ker(\chi_x|_A)=\mathfrak p_y\), and the induced absolute value on
\(\operatorname{Frac}(A/\mathfrak p_y)\) is the one determined by \(y\).
In the field \(\mathcal H(x)\) the relation
\(\chi_x(g)\chi_x(T_i)=\chi_x(f_i)\) holds with \(\chi_x(g)\neq0\), so
\(\chi_x(T_i)=\chi_x(f_i)/\chi_x(g)\) lies in the subfield
\(\operatorname{Frac}(\chi_x(A))\).  Consequently, for every polynomial
\(P\in A[T_1,\ldots,T_m]\), the value \(|P(T)(x)|\) equals the absolute
value of the corresponding element of
\(\operatorname{Frac}(A/\mathfrak p_y)\), which is determined by \(y\)
alone.  The image of \(A[T_1,\ldots,T_m]\) is dense in \(A_U\), since
polynomials are dense in the Tate algebra and quotient maps preserve
density.  Two bounded seminorms agreeing on a dense subalgebra agree,
because
\(\bigl||h|_x-|h'|_x\bigr|\leq|h-h'|_x\leq\|h-h'\|\) for the ultrametric
seminorm.  Hence \(x\) is unique over \(y\).

\emph{Homeomorphism.}
The map \(\pi\) is a continuous bijection from the compact space
\(\mathcal M(A_U)\) (\Cref{prop:berk-compact-spectrum}) onto \(U\), which
is Hausdorff as a subspace of \(\mathcal M(A)\).  A continuous bijection
from a compact space to a Hausdorff space is a homeomorphism.  Finally
\(U\) is defined by the weak inequalities \(|f_i(x)|\leq|g(x)|\) between
continuous functions of \(x\), hence is closed in \(\mathcal M(A)\), and
compact.
\end{proof}

\begin{example}[Weierstrass and Laurent domains]
\label{ex:berk-weierstrass-laurent}
Two special cases occur constantly.  Taking \(g=1\) gives the
\emph{Weierstrass domain}\index{Weierstrass domain}
\[
 \{x:|f_i(x)|\leq1,\ i=1,\ldots,m\},
 \qquad\text{represented by}\qquad
 A\langle T_1,\ldots,T_m\rangle/(T_i-f_i).
\]
Combining a Weierstrass condition with inequalities
\(|g_j(x)|\geq1\) gives a \emph{Laurent domain}\index{Laurent domain},
represented by adjoining inverse variables \(S_j\) with relations
\(g_jS_j=1\); the argument of \Cref{prop:berk-rational-domain} applies
verbatim, the character sending \(S_j\) to \(g_j(y)^{-1}\).  For instance,
inside the closed unit disc \(\mathcal M(K\langle T\rangle)\) over a
discretely valued field with uniformizer \(\pi\), the closed annulus
\[
 \{x:|\pi|\leq|T(x)|\leq1\}
 =\mathcal M\bigl(K\langle T,S\rangle/(TS-\pi)\bigr)
\]
is a Laurent domain: the relation forces \(|T(x)||S(x)|=|\pi|\) with both
factors at most \(1\), which is exactly the displayed inequality.  This
annulus reappears as the local model of a node in
\cref{sec:berk-line,chap:skeleta}.
\end{example}

\begin{remark}[Topology versus analytic domains]
A rational domain is compact, but need not be open in the ordinary topology.
It is nevertheless an analytic domain. This distinction is indispensable
for the structure sheaf, but most of pluripotential theory uses the ordinary
Hausdorff topology for upper semicontinuity, compactness, and measures.
\end{remark}

An affinoid \(K\)-analytic space is, by definition, a space locally modeled
on spectra \(\mathcal M(A)\), with the above sheaves. General Berkovich
analytic spaces are obtained by gluing affinoids along analytic domains. We
will use this construction, but no later proof depends on the full list of
sheaf axioms.

\section{Analytification of schemes}
\label{sec:berk-analytification}

Let \(X=\operatorname{Spec}A\) be an affine scheme of finite type over \(K\).
No Banach norm is chosen on \(A\).

\begin{definition}[Affine analytification]
The set \(X^{\mathrm{an}}\) consists of all multiplicative seminorms on \(A\)
extending the absolute value of \(K\). It has the topology of pointwise
convergence: the maps \(x\mapsto |f(x)|\), \(f\in A\), are required to be
continuous.
\end{definition}

The affine definition admits a useful scheme-theoretic reformulation.

\begin{proposition}\label{prop:berk-pair-description}
There is a natural bijection between \(X^{\mathrm{an}}\) and pairs
\((\xi,|\cdot|_\xi)\), where \(\xi\in X\) and \(|\cdot|_\xi\) is an absolute
value on the residue field \(\kappa(\xi)\) extending that of \(K\).
\end{proposition}

\begin{proof}
Given a seminorm \(x\), its kernel \(\mathfrak p_x\) is prime and hence
defines \(\xi\in X\). The seminorm descends to \(A/\mathfrak p_x\), is
nonzero on every nonzero element, and therefore extends uniquely to
\(\operatorname{Frac}(A/\mathfrak p_x)=\kappa(\xi)\).

Conversely, compose \(A\to\kappa(\xi)\) with the chosen absolute value.
The two operations are inverse to one another.
\end{proof}

Gluing along Zariski open subsets requires knowing that the
analytification of a basic open subscheme sits inside the affine
analytification as an open subspace.  This is the content of the next
lemma.

\begin{lemma}\label{lem:berk-localization}
Let \(X=\Spec A\) be affine of finite type over \(K\) and let \(f\in A\).
Restriction of seminorms defines a homeomorphism
\[
 (\Spec A_f)^{\mathrm{an}}
 \xrightarrow{\ \sim\ }
 \{x\in X^{\mathrm{an}}:|f(x)|\neq0\},
\]
and the target is open in \(X^{\mathrm{an}}\).
\end{lemma}

\begin{proof}
The target is the preimage of the open set
\(\mathbf R_{>0}\) under the continuous map \(x\mapsto|f(x)|\), hence
open.

A point \(y\) of \((\Spec A_f)^{\mathrm{an}}\) restricts along
\(A\to A_f\) to a multiplicative seminorm \(x\) on \(A\); since
\(f\) becomes invertible, \(|f(x)|=|f(y)|\neq0\).  Conversely, let
\(x\in X^{\mathrm{an}}\) with \(|f(x)|\neq0\).  Any extension \(y\) of
\(x\) to \(A_f\) must satisfy
\[
 |(a/f^m)(y)|=|a(x)|\,|f(x)|^{-m},
\]
by multiplicativity, so an extension is unique if it exists.  For
existence, take the displayed formula as a definition.  It is well
defined: if \(a/f^m=b/f^n\) in \(A_f\), then \(f^N(af^n-bf^m)=0\) in
\(A\) for some \(N\), so
\(|f(x)|^N|a(x)||f(x)|^n=|f(x)|^N|b(x)||f(x)|^m\); dividing by
\(|f(x)|^{N}\neq0\) gives the required equality.  Multiplicativity is
clear.  For the ultrametric inequality, bring two fractions to a common
denominator:
\[
 \Bigl|\frac{a}{f^m}+\frac{b}{f^n}\Bigr|_y
 =\frac{|af^n+bf^m|_x}{|f(x)|^{m+n}}
 \leq\max\Bigl\{\frac{|a(x)|}{|f(x)|^m},
                \frac{|b(x)|}{|f(x)|^n}\Bigr\}.
\]
Thus restriction is a bijection onto the displayed set.  The topology of
\((\Spec A_f)^{\mathrm{an}}\) is generated by the evaluations at elements
\(a/f^m\), whose absolute values are quotients of continuous functions of
\(x\) with nonvanishing denominator; conversely, evaluations at elements
of \(A\) are among the evaluations at elements of \(A_f\).  The two
topologies therefore coincide, and the bijection is a homeomorphism.
\end{proof}

For a general finite-type \(K\)-scheme, affine analytifications glue along
Zariski open subsets: \Cref{lem:berk-localization} identifies the
analytification of a principal open subscheme with an open subspace, the
identifications are compatible on triple overlaps because all of them are
restriction maps of seminorms, and principal opens form a basis. This
gives a functor
\[
              X\longmapsto X^{\mathrm{an}}.
\]
There is a continuous support map
\[
 \rho_X:X^{\mathrm{an}}\longrightarrow X,
 \qquad
 \rho_X(x)=\{f:|f(x)|=0\}
\]
when \(X\) is viewed with its Zariski topology. Its fiber above \(\xi\) is
the set of absolute values on \(\kappa(\xi)\) that extend the one on \(K\).
For a \(K\)-rational point, the given absolute value supplies a canonical
analytic point.

\begin{proposition}
\label{prop:berk-support-surjective}
For every finite-type \(K\)-scheme \(X\), the map
\(\rho_X:X^{\mathrm{an}}\to X\) is surjective.
\end{proposition}

\begin{proof}
By \Cref{prop:berk-pair-description}, surjectivity means that every
residue field \(\kappa(\xi)\), \(\xi\in X\), carries an absolute value
extending that of \(K\).  The field \(\kappa(\xi)\) is a finitely
generated extension of \(K\).  Choose a transcendence basis
\(t_1,\ldots,t_d\), so that \(\kappa(\xi)\) is a finite extension of
\(F=K(t_1,\ldots,t_d)\).

First extend the absolute value to \(F\) by the iterated Gauss norm: on
\(K[t_1]\) put \(|\sum a_it_1^i|=\max|a_i|\); this is multiplicative by
the graded-domain argument of \Cref{ex:berk-tate-disc} with \(r=1\), and
extends to the fraction field \(K(t_1)\) by multiplicativity.  Iterating
in the variables \(t_2,\ldots,t_d\) extends it to \(F\).

Next extend from \(F\) to the finite extension \(\kappa(\xi)\).  By
Chevalley's extension theorem
\cite[Chapter~VI]{ZariskiSamuelII}, the valuation
\(v=-\log|\cdot|\) of \(F\) extends to a valuation \(w\) of
\(\kappa(\xi)\) with values in some totally ordered abelian group
\(\Gamma\supset v(F^\times)\).  Because \(\kappa(\xi)/F\) is algebraic,
the quotient \(\Gamma_w/v(F^\times)\) is a torsion group: for
\(h\in\kappa(\xi)^\times\) with minimal equation
\(h^n+c_{n-1}h^{n-1}+\cdots+c_0=0\), \(c_0\ne0\), two of the terms
\(c_ih^i\) must share the minimal value---otherwise the unique term of
strictly minimal value could not cancel against the others, exactly as in
\cref{lem:strong-triangle-equality}.  Equating the values of two such
terms gives \((i-j)\,w(h)=v(c_j)-v(c_i)\in v(F^\times)\) with
\(i\neq j\), so some nonzero multiple \(kw(h)\) lies in
\(v(F^\times)\).  A torsion extension of a subgroup of \(\mathbf R\)
order-embeds into \(\mathbf R\): send \(\gamma\) with
\(n\gamma\in v(F^\times)\subset\mathbf R\) to \(n^{-1}(n\gamma)\); this
is well defined, order preserving, and additive.  Hence \(w\) may be
taken real-valued, and \(\exp(-w)\) is an absolute value on
\(\kappa(\xi)\) extending the given one on \(K\).
\end{proof}

The next lemma explains how the noncompact space \(X^{\mathrm{an}}\) is
exhausted by the compact spectra of \Cref{sec:berk-spectrum}.  It also
completes the identification promised in \Cref{ex:berk-tate-disc}.

\begin{lemma}\label{lem:berk-polydisc-exhaustion}
Let \(X=\Spec A\) with \(A=K[S_1,\ldots,S_n]/I\), and for \(r>0\) put
\[
 X^{\mathrm{an}}_{\leq r}
 =\{x\in X^{\mathrm{an}}:|S_i(x)|\leq r,\ i=1,\ldots,n\}.
\]
Then restriction of seminorms defines a homeomorphism
\[
 \mathcal M\bigl(K\langle r^{-1}S_1,\ldots,r^{-1}S_n\rangle/
 I\cdot K\langle r^{-1}S\rangle\bigr)
 \xrightarrow{\ \sim\ }
 X^{\mathrm{an}}_{\leq r}.
\]
Consequently \(X^{\mathrm{an}}=\bigcup_{r>0}X^{\mathrm{an}}_{\leq r}\) is
an increasing union of compact subsets, every point of
\(X^{\mathrm{an}}_{\leq r}\) is interior to \(X^{\mathrm{an}}_{\leq r'}\)
for \(r'>r\), and \(X^{\mathrm{an}}\) is locally compact and Hausdorff.
In particular \(\mathbf A^{1,\mathrm{an}}_K\) is the increasing union of
the closed discs \(\mathcal M(K\langle r^{-1}T\rangle)\).
\end{lemma}

\begin{proof}
Write \(B_r=K\langle r^{-1}S\rangle/I\cdot K\langle r^{-1}S\rangle\); the
ideal is closed and \(B_r\) is a generalized \(K\)-affinoid algebra,
allowing the arbitrary polyradius \(r\), with the quotient
norm \cite[\S2.1]{Berk12}.  A point of \(\mathcal M(B_r)\) is the same
thing as a bounded multiplicative seminorm on
\(K\langle r^{-1}S\rangle\) vanishing on \(I\); restricting to the dense
subring \(K[S]\) and passing to the quotient \(A\) gives a point of
\(X^{\mathrm{an}}\) with \(|S_i(x)|\leq\|S_i\|_r=r\).

Conversely, let \(x\in X^{\mathrm{an}}_{\leq r}\), viewed as a seminorm
on \(K[S]\) vanishing on \(I\).  For a polynomial
\(F=\sum_\beta a_\beta S^\beta\), the ultrametric inequality and
multiplicativity give
\[
 |F(x)|\leq\max_\beta|a_\beta|\,r^{|\beta|}=\|F\|_r,
\]
so \(x\) is norm-bounded on \(K[S]\).  Given
\(F\in K\langle r^{-1}S\rangle\) with polynomial truncations \(F_N\),
the real sequence \(|F_N(x)|\) satisfies
\(\bigl||F_N(x)|-|F_M(x)|\bigr|\leq|F_N-F_M|_x\leq\|F_N-F_M\|_r\to0\), so
it converges; define \(|F(x)|\) as the limit.  The limit does not depend
on the chosen approximating sequence, again by the displayed bound.
Multiplicativity passes to the limit because
\(\|FG-F_NG_N\|_r\to0\) and products of convergent real sequences
converge to the product of the limits; the ultrametric inequality and the
values on \(K\) pass to the limit for the same reason.  The extended
seminorm vanishes on \(I\cdot K\langle r^{-1}S\rangle\): it vanishes on
finite combinations \(\sum g_jh_j\) with \(g_j\in I\) by the ultrametric
inequality and multiplicativity, hence on their closure by continuity.
This produces the inverse map, so restriction is a bijection.

Both topologies are initial topologies: for
\(\mathcal M(B_r)\) with respect to all \(|F(\cdot)|\),
\(F\in K\langle r^{-1}S\rangle\), and for
\(X^{\mathrm{an}}_{\leq r}\) with respect to polynomials only.  The
identity map from the first to the second is continuous.  For the
converse, each \(|F(\cdot)|\) is the uniform limit on
\(\mathcal M(B_r)\) of \(|F_N(\cdot)|\), since
\(\bigl||F(x)|-|F_N(x)|\bigr|\leq\|F-F_N\|_r\); a uniform limit of
continuous functions is continuous, so the polynomial evaluations already
generate the finer topology.  Hence the bijection is a homeomorphism, and
compactness follows from \Cref{prop:berk-compact-spectrum}.

The exhaustion statement is clear since every point has
\(\max_i|S_i(x)|<\infty\).  If \(x\in X^{\mathrm{an}}_{\leq r}\) and
\(r'>r\), then the open set \(\{|S_i|<r'\}\) contains \(x\) and is
contained in \(X^{\mathrm{an}}_{\leq r'}\); thus compact neighborhoods
exist, and Hausdorffness is inherited from the seminorm description as in
\Cref{prop:berk-compact-spectrum}.
\end{proof}

We can now state the universal property of analytification.  The affine
case is a complete triviality once the definitions are in place, and we
prove it; the general case involves the structure sheaves and is cited.

\begin{proposition}
\label{prop:berk-universal-affine}
Let \(X=\Spec A\) be affine of finite type over \(K\).  For every
\(K\)-affinoid algebra \(B\) and every \(K\)-algebra homomorphism
\(\varphi:A\to B\), there is a unique continuous map
\[
 \varphi^{\mathrm{an}}:\mathcal M(B)\longrightarrow X^{\mathrm{an}}
 \qquad\text{with}\qquad
 |f(\varphi^{\mathrm{an}}(y))|=|\varphi(f)(y)|
 \quad(f\in A,\ y\in\mathcal M(B)).
\]
The assignment \(\varphi\mapsto\varphi^{\mathrm{an}}\) respects identities
and compositions.
\end{proposition}

\begin{proof}
Uniqueness is immediate: the displayed identity prescribes every value of
the seminorm \(\varphi^{\mathrm{an}}(y)\), and a point of
\(X^{\mathrm{an}}\) is nothing but its family of values.  For existence,
set \(\varphi^{\mathrm{an}}(y)=|\cdot|_y\circ\varphi\).  Each axiom of a
multiplicative seminorm extending \(|\cdot|_K\) is inherited directly
from \(y\), because \(\varphi\) is a \(K\)-algebra homomorphism.
Continuity holds because the composite of
\(y\mapsto\varphi^{\mathrm{an}}(y)\) with each generating evaluation
\(x\mapsto|f(x)|\) is the continuous map \(y\mapsto|\varphi(f)(y)|\).
Functoriality is a direct computation with composites.
\end{proof}

\begin{theorem}
\label{thm:berk-universal-property}
Let \(X\) be a scheme locally of finite type over \(K\).  There exist a
\(K\)-analytic space \(X^{\mathrm{an}}\) and a morphism of locally ringed
spaces \(\pi:X^{\mathrm{an}}\to X\) with the following universal
property: every morphism \(Y\to X\) from a good \(K\)-analytic space
\(Y\) factors uniquely through \(\pi\).  The underlying topological space
of \(X^{\mathrm{an}}\) is the glued seminorm space constructed above, and
\(\pi\) is the support map \(\rho_X\).
\end{theorem}

\begin{proof}[Foundational analytic-space input]
This is Berkovich's theorem, and we use it as stated.  What this book
actually needs is the topological content, which is already established:
the point-set space, its local compactness
(\Cref{lem:berk-polydisc-exhaustion}), the continuous support map, and
the affine factorization of \Cref{prop:berk-universal-affine}, which is
the universal property evaluated on affinoid building blocks.  The
verification that the structure sheaves assemble compatibly is
\cite[Theorem~3.4.1]{Berk12}.
\end{proof}

\begin{theorem}\label{thm:berk-proper-compact}
Let \(X\) be a finite-type \(K\)-scheme.
\begin{enumerate}
\item If \(X\) is separated, then \(X^{\mathrm{an}}\) is Hausdorff.
\item If \(X\) is proper, then \(X^{\mathrm{an}}\) is compact.
\item More precisely, analytification carries proper morphisms of finite-type
schemes to proper morphisms of analytic spaces.
\end{enumerate}
\end{theorem}
\begin{proof}[Foundational analytification input]
This is \cite[Theorem~3.4.8]{Berk12}.  That theorem states that
analytification preserves separated and proper morphisms; its properness
conclusion includes compactness after analytic base change and the absence
of relative boundary.  Apply it to the structure morphism
\(X\to\operatorname{Spec}K\) to obtain (1) and (2), and to an arbitrary
proper morphism to obtain (3).
\end{proof}

\begin{remark}
When \(K\) is algebraically closed and nontrivially valued, classical points
are dense in \(X^{\mathrm{an}}\) for a finite-type \(K\)-scheme \(X\); this
is the density of rigid points on each member of the polydisc exhaustion
\cite[Proposition~2.1.15]{Berk12}. They
are usually far from being all the points. The added nonclassical points are
what connect discs and carry model-theoretic measures.  On the affine line
the phenomenon is visible directly: by the classification of
\cref{sec:berk-line}, a Gauss point \(\zeta_{a,r}\) is the supremum of the
evaluation seminorms over the classical points of its disc, and every
basic neighborhood of \(\zeta_{a,r}\) prescribes finitely many values
\(|f_j(x)|\), which classical points of the disc realize up to any
\(\varepsilon\).  Over a trivially valued field the classical points are
the \(k\)-points and are generally not dense; see
\cref{sec:berk-trivial}.
\end{remark}

\section{Centers and reduction maps}
\label{sec:berk-reduction}

Write
\[
 K^\circ=\{a\in K:|a|\leq1\},\qquad
 K^{\circ\circ}=\{a\in K:|a|<1\},\qquad
 \widetilde K=K^\circ/K^{\circ\circ}.
\]
Let \(X\) be a proper \(K\)-scheme and let \(\mathcal X\) be a proper flat
\(K^\circ\)-model of \(X\). For \(x\in X^{\mathrm{an}}\), the absolute value
on \(\kappa(\rho_X(x))\) has valuation ring
\[
 R_x=\{h\in\kappa(\rho_X(x)):|h(x)|\leq1\}.
\]
The morphism \(\operatorname{Spec}\kappa(\rho_X(x))\to X\) extends uniquely
to \(\operatorname{Spec}R_x\to\mathcal X\) by the valuative criterion of
properness.

\begin{definition}[Center and reduction]
The image in \(\mathcal X_s\) of the closed point of
\(\operatorname{Spec}R_x\) is the \emph{center} of \(x\) on \(\mathcal X\),
denoted \(c_{\mathcal X}(x)\). The map
\[
             \operatorname{red}_{\mathcal X}:X^{\mathrm{an}}
             \longrightarrow\mathcal X_s,
             \qquad x\longmapsto c_{\mathcal X}(x),
\]
is the \emph{reduction map}.
\end{definition}

On an affine open \(\mathcal U=\operatorname{Spec}B\subset\mathcal X\), the
construction is concrete. A point has center in \(\mathcal U\) precisely
when \(|f(x)|\leq1\) for every \(f\in B\); its center then corresponds to
\[
                   \{f\in B:|f(x)|<1\}.
\]
This description also shows the characteristic reversal of topologies.

\begin{proposition}
\label{prop:berk-reduction-anticontinuous}
The reduction map is anticontinuous: inverse images of Zariski open subsets
are closed, and inverse images of Zariski closed subsets are open, relative
to the appropriate formal generic fiber.
\end{proposition}

\begin{proof}
We first make the chart-relative statement precise and then prove it.
Let \(\mathcal U=\Spec B\subset\mathcal X\) be an affine open, and set
\[
 X^{\mathrm{an}}_{\mathcal U}
 =\{x\in X^{\mathrm{an}}:
   c_{\mathcal X}(x)\in\mathcal U_s\},
\]
the generic fiber of the formal completion of \(\mathcal U\).  We claim
\[
 X^{\mathrm{an}}_{\mathcal U}
 =\{x\in X^{\mathrm{an}}:\rho_X(x)\in\mathcal U_K
   \text{ and }|f(x)|\leq1\ \text{for all }f\in B\},
\]
where \(\mathcal U_K=\mathcal U\cap X\).  Indeed, if
\(c_{\mathcal X}(x)\in\mathcal U\), then the support
\(\xi=\rho_X(x)\), being a generization of the center, also lies in the
open set \(\mathcal U\); and the composite
\(B\to\mathcal O_{\mathcal X,c(x)}\to R_x\) shows \(|f(x)|\leq1\) for
\(f\in B\).  Conversely, if \(\xi\in\mathcal U_K\) and every \(f\in B\)
satisfies \(|f(x)|\leq1\), then \(B\to\kappa(\xi)\) lands in \(R_x\), so
\(\operatorname{Spec}R_x\to\mathcal U\subset\mathcal X\) extends the
generic-point morphism; by separatedness this is the unique extension, and
its closed point lies in \(\mathcal U_s\).  The set
\(X^{\mathrm{an}}_{\mathcal U}\) is compact: choosing finitely many
\(K^\circ\)-algebra generators \(b_1,\ldots,b_N\) of \(B\), it is the
subset \(\{|b_j(x)|\leq1,\ j=1,\ldots,N\}\) of
\((\Spec B_K)^{\mathrm{an}}\), which is compact by
\Cref{lem:berk-polydisc-exhaustion} applied to the presentation by the
\(b_j\) with \(r=1\).

Now fix \(\tilde f\in\widetilde B=B\otimes_{K^\circ}\widetilde K\) and a
lift \(f\in B\).  For \(x\in X^{\mathrm{an}}_{\mathcal U}\), the center
corresponds to the prime \(\{g\in B:|g(x)|<1\}\), so
\[
 \operatorname{red}_{\mathcal X}(x)\in V(\tilde f)
 \iff|f(x)|<1,
 \qquad
 \operatorname{red}_{\mathcal X}(x)\in D(\tilde f)
 \iff|f(x)|=1 .
\]
The condition is independent of the lift: two lifts differ by an element
\(g=\sum_ja_jb_j'\) with \(a_j\in K^{\circ\circ}\) and \(b_j'\in B\), so
\(|g(x)|\leq\max_j|a_j|<1\), and \cref{lem:strong-triangle-equality}
shows that the conditions \(|f(x)|<1\) and \(|f(x)|=1\) are unaffected.
On \(X^{\mathrm{an}}_{\mathcal U}\), the first displayed condition is a
strict inequality between continuous functions, hence defines a
relatively open subset; the second is the weak inequality
\(|f(x)|\geq1\) inside the region \(|f(x)|\leq1\), hence defines a
relatively closed subset.

The special fiber \(\mathcal U_s=\Spec\widetilde B\) is noetherian
because \(\widetilde B\) is a finite-type \(\widetilde K\)-algebra.  A
Zariski closed subset is \(V(\tilde f_1,\ldots,\tilde f_k)
=\bigcap_jV(\tilde f_j)\), whose inverse image is a finite intersection
of relatively open sets, hence relatively open.  A Zariski open subset is
the complement of a closed one; since
\(\operatorname{red}_{\mathcal X}^{-1}\) turns complements in
\(\mathcal U_s\) into complements in \(X^{\mathrm{an}}_{\mathcal U}\),
its inverse image is relatively closed.  This proves the chart-relative
statement.

Finally, for a proper model the assertion globalizes.  Choose a finite
affine cover \(\mathcal U_1,\ldots,\mathcal U_m\) of \(\mathcal X\); the
compact sets \(X^{\mathrm{an}}_{\mathcal U_i}\) cover
\(X^{\mathrm{an}}\), since every point has a center.  If
\(W\subset\mathcal X_s\) is Zariski open, then
\(\operatorname{red}_{\mathcal X}^{-1}(W)\cap
X^{\mathrm{an}}_{\mathcal U_i}\) is closed in the compact set
\(X^{\mathrm{an}}_{\mathcal U_i}\), hence compact; the union of these
finitely many compact sets is compact, hence closed in the Hausdorff
space \(X^{\mathrm{an}}\).  Taking complements shows that inverse images
of Zariski closed subsets of \(\mathcal X_s\) are open in
\(X^{\mathrm{an}}\).
\end{proof}

Assume now that \(K\) is discretely valued, choose a uniformizer \(\pi\), and
normalize \(v(\pi)=1\). Let \(\mathcal X\) be normal and write
\[
                        \mathcal X_s=\sum_E b_EE.
\]
At the generic point of \(E\), the local ring is a discrete valuation ring
and \(\operatorname{ord}_E(\pi)=b_E\). Hence
\[
                  v_E=\frac{1}{b_E}\operatorname{ord}_E
\]
extends \(v\) on \(K\). It defines a point \(x_E\in X^{\mathrm{an}}\), called
the \emph{divisorial point} associated to \(E\). Its reduction is the
generic point of \(E\). Normality implies that it is the unique analytic
point reducing to that generic point.  Because the model Monge--Amp\`ere
measures of \cref{chap:metrics} are supported at divisorial points, this
uniqueness deserves a complete proof.

\begin{proposition}
\label{prop:berk-divisorial-fibre}
Let \(K\) be discretely valued, let \(X\) be an integral proper
\(K\)-scheme, and let \(\mathcal X\) be a normal proper flat model.  For
every irreducible component \(E\) of \(\mathcal X_s\), with generic point
\(\xi_E\),
\[
 \operatorname{red}_{\mathcal X}^{-1}(\xi_E)=\{x_E\}.
\]
\end{proposition}

\begin{proof}
That \(\operatorname{red}_{\mathcal X}(x_E)=\xi_E\) is
\Cref{prop:component-valuation}: the valuation ring of \(v_E\) dominates
the local ring at \(\xi_E\).

Conversely, let \(x\in X^{\mathrm{an}}\) with
\(\operatorname{red}_{\mathcal X}(x)=\xi_E\), and write
\(\xi=\rho_X(x)\in X\) for its support and
\(\mathcal O=\mathcal O_{\mathcal X,\xi_E}\).

\emph{The support is the generic point of \(X\).}
Since \(X\) is integral and \(\mathcal X\) is a normal model with
irreducible generic fiber, \(\mathcal X\) is integral: flatness forces
every irreducible component of \(\mathcal X\) to meet the generic fiber
in a nonempty open subset of the irreducible space \(X\), and normality
of the connected components then leaves a single component.  The point
\(\xi_E\) is a minimal prime over the principal ideal \((\pi)\), so Krull's
principal ideal theorem gives \(\dim\mathcal O\le 1\); as \(\xi_E\) is not
the generic point of \(\mathcal X\), we get \(\dim\mathcal O=1\).  The
center \(\xi_E\) is a specialization of \(\xi\), so \(\xi\) corresponds
to a prime \(\mathfrak q\) of the one-dimensional local ring
\(\mathcal O\), namely \(\mathfrak q=0\) or the maximal ideal.  Since
\(\xi\) lies in the generic fiber, \(\pi\notin\mathfrak q\), while
\(\pi\) lies in the maximal ideal; hence \(\mathfrak q=0\) and \(\xi\)
is the generic point of \(\mathcal X\), that is, of \(X\).  In
particular \(x\) has support zero and is a genuine valuation on
\(K(X)\).

\emph{Rigidity of the dominated ring.}
The ring \(\mathcal O\) is a one-dimensional normal noetherian local
domain, hence a discrete valuation ring with fraction field \(K(X)\).
The valuation ring \(R_x\subset K(X)\) dominates \(\mathcal O\) by the
definition of the center.  We claim that any ring \(R\) with
\(\mathcal O\subseteq R\subseteq K(X)\) equals \(\mathcal O\) or
\(K(X)\).  Let \(u\in R\setminus\mathcal O\), and let
\(\varpi\in\mathcal O\) be a uniformizer of \(\mathcal O\).  Then
\(u=w\varpi^{-k}\) with \(w\in\mathcal O^\times\) and \(k\geq1\), so
\(\varpi^{-k}=uw^{-1}\in R\) and
\(\varpi^{-1}=\varpi^{k-1}\varpi^{-k}\in R\); therefore
\(R\supseteq\mathcal O[\varpi^{-1}]=K(X)\).  Since \(R_x\) dominates
\(\mathcal O\), its maximal ideal meets \(\mathcal O\) in the maximal
ideal, so \(R_x\neq K(X)\); hence \(R_x=\mathcal O\).

\emph{Normalization.}
Equality of valuation rings inside \(K(X)\) with value groups in
\(\mathbf R\) means \(v_x=c\operatorname{ord}_E\) for some \(c>0\).
Evaluating at \(\pi\) gives
\(1=v_x(\pi)=c\operatorname{ord}_E(\pi)=cb_E\), so \(c=b_E^{-1}\) and
\(v_x=v_E\), that is, \(x=x_E\).
\end{proof}

\begin{remark}[Surjectivity of reduction]
\label{rem:berk-reduction-surjective}
For a strictly \(K\)-affinoid algebra \(A\), the reduction map
\(\mathcal M(A)\to\Spec\widetilde A\) to the spectrum of the canonical
reduction is surjective; this is \cite[Proposition~2.4.4]{Berk12} and
uses the theory of the spectral seminorm.  Combined with the valuative
criterion, it implies that the reduction map of a proper flat model hits
every closed point of the special fiber.  We use surjectivity only
through the divisorial points, whose existence was proved by hand above,
so the deep input is quarantined in this remark.
\end{remark}

Divisorial points will support the model Monge--Amp\`ere measures of
\cref{chap:metrics}. The passage from a component to \(x_E\) explains the factor
\(b_E^{-1}\) in the valuation and the factor \(b_E\) in the corresponding
measure.  We close the section with two fully computed reduction maps;
the second also previews the skeleta of \cref{chap:skeleta}.

\begin{example}[Reduction on \(\mathbf P^1\) over \(\mathbf Q_p\)]
\label{ex:berk-reduction-p1-qp}
Let \(K=\mathbf Q_p\) with \(|p|=p^{-1}\), let
\(X=\mathbf P^1_{\mathbf Q_p}\), and let
\(\mathcal X=\mathbf P^1_{\mathbf Z_p}\), covered by the two affine
charts \(\Spec\mathbf Z_p[T]\) and \(\Spec\mathbf Z_p[T^{-1}]\).  The
special fiber is \(\mathbf P^1_{\mathbf F_p}\), reduced and irreducible,
with \(b=1\).

\emph{Points with \(|T(x)|\leq1\)} have center in the first chart, and
the center is the prime \(\{f\in\mathbf Z_p[T]:|f(x)|<1\}\).
\begin{enumerate}
\item For the evaluation point at \(a\in\mathbf Z_p\), one has
  \(|f(x)|=|f(a)|<1\) exactly when \(f(a)\in p\mathbf Z_p\), i.e.\ when
  \(\tilde f(\tilde a)=0\).  The center is the closed point
  \(\tilde a\in\mathbf A^1_{\mathbf F_p}\).
\item For the Gauss point \(\zeta_{0,1}\),
  \(|\sum a_iT^i|=\max|a_i|<1\) means that all coefficients lie in
  \(p\mathbf Z_p\), i.e.\ \(\tilde f=0\).  The center is the generic
  point of the special fiber, as
  \Cref{prop:berk-divisorial-fibre} predicts: \(\zeta_{0,1}\) is the
  divisorial point of the unique component.
\item For \(\zeta_{a,r}\) with \(a\in\mathbf Z_p\) and \(0<r<1\),
  compute with \(f=\sum c_i(T-a)^i\), \(c_i\in\mathbf Z_p\):
  \(|f(\zeta_{a,r})|=\max_i|c_i|r^i\).  For \(i\ge1\) each term is at
  most \(r<1\), so \(|f(\zeta_{a,r})|<1\) if and only if \(|c_0|<1\),
  i.e.\ \(\tilde f(\tilde a)=0\).  The center is again \(\tilde a\).
\item For a point with support at a nonrational closed point, such as the
  evaluation at \(\alpha=\sqrt p\) (the seminorm
  \(f\mapsto|f(\alpha)|\) with kernel \((T^2-p)\)), one has
  \(|\alpha|=p^{-1/2}<1\), so \(|f(\alpha)|<1\) exactly when the reduction
  of \(f\) vanishes at \(\tilde\alpha=0\); the center is
  \(\tilde 0\).  Closed points of higher residue degree occur as well:
  if \(q\in\mathbf Z_p[T]\) is monic with irreducible reduction
  \(\tilde q\) of degree \(d\ge2\), the evaluation points at the roots
  of \(q\) reduce to the degree-\(d\) closed point \(\{\tilde q=0\}\)
  of \(\mathbf A^1_{\mathbf F_p}\).
\end{enumerate}
\emph{Points with \(|T(x)|>1\)} have center in the second chart at the
point \(\widetilde\infty\), since \(|T^{-1}(x)|<1\).

Assembling these computations,
\[
 \operatorname{red}^{-1}(\text{generic point})=\{\zeta_{0,1}\},
 \qquad
 \operatorname{red}^{-1}(\tilde a)
 =\{x:|(T-a)(x)|<1\}
 \quad(a\in\mathbf Z_p),
\]
and similarly for \(\widetilde\infty\) and for closed points of higher
degree.  The fiber over each closed point is an open disc-like domain,
open in \(X^{\mathrm{an}}\), while the fiber over the generic point is the
single Gauss point, and its complement is the disjoint union of the open
fibers: anticontinuity (\Cref{prop:berk-reduction-anticontinuous}) is
visible in every clause.
\end{example}

\begin{example}[Reduction to a two-component model over
\(\mathbf C(\!(t)\!)\)]
\label{ex:berk-reduction-p1-cct}
Let \(K=\mathbf C(\!(t)\!)\) with \(v(t)=1\), so
\(K^\circ=\mathbf C[\![t]\!]\) and \(\widetilde K=\mathbf C\).  Let
\(X=\mathbf P^1_K\) and let
\(\mathcal X\to\mathbf P^1_{K^\circ}\) be the blowup at the origin of the
special fiber.  Around the blown-up point, \(\mathcal X\) is covered by
the two charts
\[
 \mathcal U_A=\Spec K^\circ[u],\quad T=tu,
 \qquad
 \mathcal U_B=\Spec K^\circ[T,w]/(t-Tw),
\]
with \(u=T/t\) and \(w=t/T=u^{-1}\) on overlaps.  The special fiber
\(\{t=0\}\) has two components, both reduced: the exceptional line \(E\),
with coordinate \(u\) (equivalently \(1/w\)), and the strict transform
\(\widetilde E\) of the original special fiber, with coordinate \(T\);
they meet at the node \(T=w=0\) of chart \(\mathcal U_B\).

Both components have multiplicity one, since
\(\operatorname{div}_{\mathcal X}(t)=E+\widetilde E\) by the equation
\(t=Tw\).  The divisorial points are computed from
\(\operatorname{ord}\): along \(\widetilde E\) one has
\(\operatorname{ord}_{\widetilde E}(T)=0\) and
\(\operatorname{ord}_{\widetilde E}(t)=1\), so
\(x_{\widetilde E}=\zeta_{0,1}\) is the Gauss point of the unit disc;
along \(E\) one has \(\operatorname{ord}_E(T)=1\) and
\(\operatorname{ord}_E(t)=1\), so \(x_E=\zeta_{0,|t|}\) is the Gauss
point of the disc of radius \(|t|=\rme^{-1}\).

The reduction map on the region \(|t|\leq|T(x)|\leq1\) is now completely
described by the chart criterion of
\Cref{prop:berk-reduction-anticontinuous}:
\begin{enumerate}
\item \(\operatorname{red}^{-1}(\text{generic point of }\widetilde E)
      =\{\zeta_{0,1}\}\) and
      \(\operatorname{red}^{-1}(\text{generic point of }E)
      =\{\zeta_{0,|t|}\}\), by
      \Cref{prop:berk-divisorial-fibre}.
\item The fiber over the node is the set of points centered at
      \(T=w=0\) in chart \(\mathcal U_B\), i.e.\ satisfying
      \(|T(x)|<1\) and \(|w(x)|=|t|/|T(x)|<1\):
      \[
       \operatorname{red}^{-1}(\text{node})
       =\{x:|t|<|T(x)|<1\},
      \]
      the open annulus.  In particular the whole open segment of Gauss
      points \(\zeta_{0,r}\), \(|t|<r<1\), reduces to the node.  This
      open annulus retracts onto that segment, which becomes an edge of
      the dual complex in \cref{chap:skeleta}; compare the standard
      annulus of \Cref{ex:berk-standard-annulus} via \(x=T\),
      \(y=t/T\).
\item The fiber over a closed point \(\tilde u=\tilde c\) of
      \(E\setminus\{\text{node}\}\), with \(c\in(K^\circ)^\times\), is
      \(\{x:|(T-ct)(x)|<|t|\}\), an open disc of radius \(|t|\); the
      fiber over the origin \(\tilde u=0\) of \(E\) is
      \(\{x:|T(x)|<|t|\}\).
\item The fiber over a closed point \(\tilde a\neq\tilde 0\) of
      \(\widetilde E\setminus\{\text{node}\}\), with \(|a|=1\), is
      \(\{x:|(T-a)(x)|<1\}\), exactly as on the unblown model.
\end{enumerate}
Blowing up has therefore split the old fiber over \(\tilde 0\) into the
divisorial point \(\zeta_{0,|t|}\), the annulus over the node, the discs
of radius \(|t|\) over the new closed points, and the disc
\(\{|T|<|t|\}\); refining the model refines the reduction.
\end{example}

\section{The analytic affine and projective lines}
\label{sec:berk-line}

The set \(\mathbf A^{1,\mathrm{an}}_K\) consists of multiplicative seminorms
on \(K[T]\). For \(a\in K\) and \(r\geq0\), define
\[
 \left|\sum_{i=0}^d c_i(T-a)^i\right|_{a,r}
       =\max_i |c_i|r^i.
\]
For \(r=0\) this is evaluation at \(a\); for \(r>0\) it is the Gauss norm of
the closed disc
\[
 D(a,r)=\{c\in K:|c-a|\leq r\}.
\]
Two discs give the same seminorm precisely when
they are the same disc.  This assertion, and everything else in this
section, rests on the elementary combinatorics of ultrametric discs,
which we record once and for all.  We write \(x\leq y\) for points of
\(\mathbf A^{1,\mathrm{an}}_K\) if \(|f(x)|\leq|f(y)|\) for every
\(f\in K[T]\).\index{Berkovich affine line!partial order}

\begin{lemma}\label{lem:berk-disc-combinatorics}
Let \(a,b\in K\) and \(r,s\geq0\).
\begin{enumerate}
\item If \(b\in D(a,r)\), then \(D(b,r)=D(a,r)\): every point of an
      ultrametric disc is a center.
\item Two discs are either disjoint or nested.
\item \(\zeta_{a,r}\leq\zeta_{b,s}\) if and only if
      \(D(a,r)\subseteq D(b,s)\), if and only if \(|a-b|\leq s\) and
      \(r\leq s\).  In particular \(\zeta_{a,r}=\zeta_{b,s}\) if and
      only if \(D(a,r)=D(b,s)\).
\item \(|(T-c)(\zeta_{a,r})|=\max\{|a-c|,r\}\) for every \(c\in K\).
\end{enumerate}
\end{lemma}

\begin{proof}
(1) If \(|c-b|\leq r\) and \(|b-a|\leq r\), then
\(|c-a|\leq\max\{|c-b|,|b-a|\}\leq r\), and symmetrically.

(2) If \(c\in D(a,r)\cap D(b,s)\) with, say, \(r\leq s\), then by (1)
\(D(a,r)=D(c,r)\subseteq D(c,s)=D(b,s)\).

(4) Expand \(T-c=(T-a)+(a-c)\); the definition of \(\zeta_{a,r}\) gives
\(|(T-c)(\zeta_{a,r})|=\max\{r,|a-c|\}\).

(3) Suppose \(|a-b|\leq s\) and \(r\leq s\).  For
\(f=\sum_ic_i(T-b)^i\), rewrite around \(a\):
\(f=\sum_jc_j'(T-a)^j\) with
\(c_j'=\sum_{i\geq j}\binom{i}{j}c_i(a-b)^{i-j}\), so
\(|c_j'|\leq\max_{i\geq j}|c_i|s^{i-j}\), using \(|a-b|\le s\) and
\(|\binom ij|\le1\).  Hence
\[
 |f(\zeta_{a,r})|=\max_j|c_j'|r^j
 \leq\max_{i\geq j}|c_i|s^{i-j}r^j
 \leq\max_i|c_i|s^i=|f(\zeta_{b,s})|,
\]
using \(r\leq s\).  Conversely, if \(\zeta_{a,r}\leq\zeta_{b,s}\), then
by (4)
\[
 \max\{|a-b|,r\}=|(T-b)(\zeta_{a,r})|\leq|(T-b)(\zeta_{b,s})|=s,
\]
which gives both inequalities.  The inclusion
\(D(a,r)\subseteq D(b,s)\) is equivalent to the same two inequalities by
(1) and (2).  For the final statement, equality of seminorms gives
mutual inclusion of discs by the above; conversely equal discs have
equal Gauss seminorms by (3) applied twice.
\end{proof}

Suppose in the rest of this section that \(K\) is algebraically closed and
nontrivially valued. Every point of the analytic affine line is of one of
the following four types.
\begin{enumerate}
\item Type I: an evaluation point \(a\in K\), equivalently \(r=0\).
\item Type II: a disc point \(\zeta_{a,r}\) with
      \(r\in |K^\times|\).
\item Type III: a disc point \(\zeta_{a,r}\) with
      \(r\notin |K^\times|\).
\item Type IV: the limit of a decreasing family of closed discs with empty
      intersection.
\end{enumerate}
Type IV points occur only when \(K\) is not spherically complete. They are
not an artificial addition: without them the space need not be complete or
compact in the relevant topology.  We now prove the classification in
full; it is the pedagogical cornerstone of this chapter.  The hypotheses
in force are: \(K\) complete, algebraically closed, and nontrivially
valued.  Completeness enters only once, algebraic closedness enters
through factorization into linear polynomials, and nontriviality is never
used; the trivially valued case is discussed in
\cref{sec:berk-trivial}.

\begin{theorem}
\label{thm:berk-a1-classification}
\index{Berkovich affine line!classification of points}
For \(x\in\mathbf A^{1,\mathrm{an}}_K\) and \(a\in K\), set
\[
 r_a(x)=|(T-a)(x)|,\qquad
 D_a(x)=D\bigl(a,r_a(x)\bigr),\qquad
 E(x)=\{D_a(x):a\in K\}.
\]
\begin{enumerate}
\item\label{itm:berk-class:nested}
  The family \(E(x)\) is totally ordered by inclusion, \(x\leq\zeta_D\)
  for every \(D\in E(x)\), and
  \(|f(x)|=\inf_{D\in E(x)}|f(\zeta_D)|\) for every \(f\in K[T]\).
\item\label{itm:berk-class:nonempty}
  If \(\bigcap E(x)\neq\varnothing\), then
  \(\bigcap E(x)=D(b,r)\) for every \(b\in\bigcap E(x)\), where
  \(r=\min_{a\in K}r_a(x)=r_b(x)\), and \(x=\zeta_{b,r}\).  Thus \(x\)
  is of type I if \(r=0\), of type II if \(r\in|K^\times|\), and of
  type III otherwise.
\item\label{itm:berk-class:empty}
  If \(\bigcap E(x)=\varnothing\), then
  \(r_\infty(x):=\inf_ar_a(x)\) is strictly positive and not attained;
  for every fixed \(f\in K[T]\), the value \(|f(\zeta_D)|\) is
  independent of \(D\in E(x)\) once \(D\) is small enough, and equals
  \(|f(x)|\).  In particular the net \((\zeta_D)_{D\in E(x)}\), indexed
  by \(E(x)\) ordered by reverse inclusion, converges to \(x\), and
  \(x\) is of type IV.
\item\label{itm:berk-class:exclusive}
  The four types are mutually exclusive, so every point of
  \(\mathbf A^{1,\mathrm{an}}_K\) is of exactly one type.
\end{enumerate}
\end{theorem}

\begin{proof}
\emph{Step 1: the one-sided bound.}
Fix \(a\in K\) and expand \(f=\sum_ic_i(T-a)^i\).  Multiplicativity
gives \(|(T-a)^i(x)|=r_a(x)^i\), so the ultrametric inequality yields
\[
 |f(x)|\leq\max_i|c_i|\,r_a(x)^i=|f(\zeta_{a,r_a(x)})|.
\]
Hence \(x\leq\zeta_D\) for every \(D\in E(x)\), and
\(|f(x)|\le\inf_{D\in E(x)}|f(\zeta_D)|\).

\emph{Step 2: nestedness.}
Let \(a,c\in K\).  If \(|a-c|\leq r_a(x)\), then
\(r_c(x)\leq\max\{r_a(x),|a-c|\}=r_a(x)\) and \(c\in D_a(x)\), so
\(D_c(x)\subseteq D(c,r_a(x))=D_a(x)\) by part~(1) of
\Cref{lem:berk-disc-combinatorics}.
If \(|a-c|>r_a(x)\), the strong triangle equality
(\cref{lem:strong-triangle-equality}) gives \(r_c(x)=|a-c|\), so
\(D_c(x)=D(c,|a-c|)\ni a\) and
\(D_a(x)\subseteq D(a,|a-c|)=D_c(x)\).  Any two members of \(E(x)\) are
therefore nested.

\emph{Step 3: the case of a common point.}
Assume \(b\in\bigcap E(x)\).  For every \(a\), membership
\(|b-a|\leq r_a(x)\) gives
\(r_b(x)\leq\max\{r_a(x),|a-b|\}=r_a(x)\); hence
\(r:=r_b(x)=\min_ar_a(x)\).  We claim \(x=\zeta_{b,r}\).  Since both
seminorms are multiplicative and every monic polynomial splits into
linear factors over the algebraically closed field \(K\), it suffices to
prove \(|(T-c)(x)|=|(T-c)(\zeta_{b,r})|\) for every \(c\in K\).  If
\(|b-c|>r\), the strong triangle equality applied to
\((T-c)=(T-b)+(b-c)\) gives \(r_c(x)=|b-c|\), while
\Cref{lem:berk-disc-combinatorics} gives
\(|(T-c)(\zeta_{b,r})|=\max\{|b-c|,r\}=|b-c|\).  If \(|b-c|\leq r\),
then \(c\in D(b,r)\subseteq D_a(x)\) for every \(a\), so \(c\) is
itself a common point; the minimality just proved, applied to \(c\),
gives \(r_c(x)=r\), while
\(|(T-c)(\zeta_{b,r})|=\max\{|b-c|,r\}=r\).  In both cases the values
agree, so \(x=\zeta_{b,r}\).  Finally
\(D(b,r)\subseteq\bigcap E(x)\subseteq D_b(x)=D(b,r)\), which
identifies the intersection.

\emph{Step 4: the case of empty intersection.}
Assume \(\bigcap E(x)=\varnothing\), and set
\(r_\infty=\inf_ar_a(x)\).

First, \(r_\infty>0\).  Otherwise choose \(a_n\) with
\(r_{a_n}(x)\to0\).  For \(m,n\), nestedness gives
\(|a_n-a_m|\leq\max\{r_{a_n}(x),r_{a_m}(x)\}\), so \((a_n)\) is Cauchy;
by completeness of \(K\) it converges to some \(a_*\in K\).  Then
\[
 r_{a_*}(x)\leq\max\{r_{a_n}(x),|a_n-a_*|\}\longrightarrow0,
\]
so \(r_{a_*}(x)=0\) and \(D_{a_*}(x)=\{a_*\}\); by nestedness
\(a_*\in D_a(x)\) for every \(a\), contradicting empty intersection.
This is the one use of completeness.

Second, the infimum is not attained.  If \(r_a(x)=r_\infty\), take any
\(c\): the discs \(D_a(x)\) and \(D_c(x)\) are nested; if
\(D_c(x)\subseteq D_a(x)\) then \(r_c(x)\ge r_\infty=r_a(x)\) forces
equal radii, hence \(D_c(x)=D_a(x)\ni a\); if
\(D_a(x)\subseteq D_c(x)\) then \(a\in D_c(x)\) directly.  Either way
\(a\in\bigcap E(x)\), a contradiction.  Consequently, for every
\(D\in E(x)\) there is \(D'\in E(x)\) with \(D'\subsetneq D\) and
strictly smaller radius.

Third, eventual constancy.  Fix \(c\in K\).  Since the intersection is
empty, there is \(a\) with \(c\notin D_a(x)\), i.e.\
\(|c-a|>r_a(x)\); then \(r_c(x)=|c-a|\) by the strong triangle
equality.  For any \(a'\) with \(D_{a'}(x)\subseteq D_a(x)\) we get
\(|a-a'|\leq r_a(x)<|c-a|\), hence \(|c-a'|=|c-a|>r_a(x)\ge
r_{a'}(x)\), and therefore
\[
 |(T-c)(\zeta_{D_{a'}(x)})|
 =\max\{|c-a'|,r_{a'}(x)\}
 =|c-a|
 =r_c(x)
 =|(T-c)(x)|.
\]
So the value of every linear polynomial along the net stabilizes at its
value at \(x\).  For a general
\(f=u\prod_{j=1}^d(T-c_j)\), choose for each \(j\) a disc
\(D_j\in E(x)\) beyond which the factor \(T-c_j\) has stabilized, and
let \(D_*\) be the smallest of the finitely many \(D_j\) (they are
nested).  For every \(D\subseteq D_*\) in \(E(x)\), multiplicativity of
Gauss seminorms and of \(x\) gives
\[
 |f(\zeta_D)|=|u|\prod_j|(T-c_j)(\zeta_D)|
            =|u|\prod_j|(T-c_j)(x)|=|f(x)|.
\]
Combined with Step 1, this proves the displayed infimum formula of
part~(1) in this case, shows that the net
\((\zeta_D)_D\) converges pointwise, hence in
\(\mathbf A^{1,\mathrm{an}}_K\), to \(x\), and exhibits \(x\) as the
limit of the strictly decreasing family \(E(x)\) with empty
intersection: \(x\) is of type IV.  In the case of Step 3 the infimum
formula holds as well, because the infimum is attained at
\(D_b(x)\) and \(x=\zeta_{D_b(x)}\).

\emph{Step 5: exclusivity.}
If \(x=\zeta_{b,s}\) is of type I, II, or III, then
\(b\in\bigcap E(x)\): indeed
\(r_a(\zeta_{b,s})=\max\{|a-b|,s\}\geq|a-b|\) for every \(a\), so
\(b\in D_a(x)\).  Hence points with \(\bigcap E(x)=\varnothing\) are of
none of the first three types, and conversely.  Within the first three
types, \(x=\zeta_{b,s}\) determines its disc: by
\Cref{lem:berk-disc-combinatorics}, \(\zeta_{b,s}=\zeta_{b',s'}\)
implies \(D(b,s)=D(b',s')\), so the radius \(s=\min_ar_a(x)\) is an
invariant of the point, and the three conditions \(s=0\),
\(s\in|K^\times|\), and \(s>0\), \(s\notin|K^\times|\) are mutually
exclusive.
\end{proof}

\begin{corollary}
\label{cor:berk-gauss-sup}
For every \(f\in K[T]\), \(a\in K\), and \(r>0\),
\[
 |f(\zeta_{a,r})|=\sup_{c\in D(a,r)}|f(c)|,
\]
and the supremum is attained whenever \(r\in|K^\times|\).  Moreover the
classical points of \(D(a,r)\) are dense in the analytic closed disc
\(\{x\in\mathbf A^{1,\mathrm{an}}_K:r_a(x)\leq r\}
=\{x:x\leq\zeta_{a,r}\}\).
\end{corollary}

\begin{proof}
The equality of the two descriptions of the analytic disc follows from
\Cref{thm:berk-a1-classification}: if \(r_a(x)\leq r\) then
\(x\leq\zeta_{a,r_a(x)}\leq\zeta_{a,r}\), and conversely
\(x\le\zeta_{a,r}\) gives
\(r_a(x)\le|(T-a)(\zeta_{a,r})|=r\).

For \(c\in D(a,r)\), the evaluation point satisfies
\(\zeta_{c,0}\leq\zeta_{a,r}\) by \Cref{lem:berk-disc-combinatorics},
which proves the inequality \(\geq\) in the display.  For \(\leq\),
factor \(f=u\prod_{j=1}^d(T-c_j)\), so that
\(|f(\zeta_{a,r})|=|u|\prod_j\max\{|a-c_j|,r\}\).  Since \(K\) is
algebraically closed, the group \(|K^\times|\) is divisible, hence dense
in \(\mathbf R_{>0}\).  Choose \(s=|\varrho|\in|K^\times|\) with
\(s\leq r\), and test points \(c=a+\varrho\theta\) with
\(\theta\in K^\circ\), \(|\theta|=1\).  For a root with \(|a-c_j|>s\),
\(|c-c_j|=|a-c_j|\) for every such \(\theta\).  For a root with
\(|a-c_j|\leq s\), put \(d_j=(c_j-a)/\varrho\in K^\circ\); then
\(|c-c_j|=s|\theta-d_j|\), which equals \(s\) as soon as the residue
\(\tilde\theta\) differs from \(\tilde d_j\).  The residue field of an
algebraically closed valued field is infinite, so a residue avoiding
the finitely many \(\tilde d_j\) exists, giving
\[
 |f(c)|=|u|\prod_j\max\{|a-c_j|,s\}.
\]
If \(r\in|K^\times|\), take \(s=r\): the supremum is attained.  In
general, letting \(s\uparrow r\) through \(|K^\times|\) makes the right
side converge to \(|f(\zeta_{a,r})|\), which proves the supremum
formula.

For density, let \(x\leq\zeta_{a,r}\) and let
\(U=\{y:\alpha_i<|f_i(y)|<\beta_i,\ i=1,\ldots,k\}\) be a basic
neighborhood of \(x\).  If \(x\) is of type IV, the net of Gauss points
of \(E(x)\) converges to \(x\) by
\Cref{thm:berk-a1-classification}, and its deep members lie below
\(\zeta_{a,r}\); replacing \(x\) by one of them, we may assume
\(x=\zeta_{b,s_0}\) with \(D(b,s_0)\subseteq D(a,r)\), the type-I case
being trivial.  Run the previous paragraph simultaneously for the
finitely many \(f_i\), with one \(\theta\) whose residue avoids the
finitely many excluded residues of all \(f_i\) at once: for
\(s\in|K^\times|\), \(s\le s_0\), we obtain classical
\(c\in D(b,s_0)\subseteq D(a,r)\) with
\(|f_i(c)|=|u_i|\prod_j\max\{|b-c_{ij}|,s\}\).  As \(s\uparrow s_0\)
these values increase to \(|f_i(x)|<\beta_i\), so for \(s\) close
enough to \(s_0\) they lie in \((\alpha_i,\beta_i)\) for every \(i\).
Then \(c\in U\).
\end{proof}

The four types can be told apart by intrinsic invariants of the
completed residue field, without reference to a family of discs.  The
statement uses the Abhyankar inequality in a version with a valued
ground field; the proof is the same independence argument as in
\Cref{thm:abhyankar-inequality}, and we spell it out.

\begin{lemma}
\label{lem:berk-abhyankar-valued}
Let \(K\) be a non-Archimedean valued field with residue field
\(\widetilde K\) and value group \(\Gamma_K=v(K^\times)\subset\mathbf R\).
Let \(F/K\) be a finitely generated field extension of transcendence
degree \(n\), and let \(w\) be a real valuation of \(F\) extending
\(v\).  Then
\[
 \dim_{\mathbf Q}
 \bigl((\Gamma_w/\Gamma_K)\otimes_{\mathbf Z}\mathbf Q\bigr)
 +\operatorname{trdeg}_{\widetilde K}\,\kappa(w)
 \;\leq\;n,
\]
where \(\Gamma_w=w(F^\times)\) and \(\kappa(w)\) is the residue field of
\(w\).
\end{lemma}

\begin{proof}
Choose \(f_1,\ldots,f_r\in F^\times\) whose values
\(\gamma_i=w(f_i)\) are \(\mathbf Q\)-linearly independent modulo
\(\Gamma_K\otimes\mathbf Q\), and units \(g_1,\ldots,g_s\in R_w^\times\)
whose residues are algebraically independent over \(\widetilde K\).  It
suffices to show that \(f_1,\ldots,f_r,g_1,\ldots,g_s\) are
algebraically independent over \(K\); the claim then follows by taking
families of maximal size, since
\(\operatorname{trdeg}_KF=n\).

Suppose there is a nonzero polynomial relation.  Group its monomials
according to the exponents of the \(f_i\):
\[
 \sum_{\alpha\in\mathbf N^r}f^\alpha P_\alpha(g)=0,
\]
with finitely many nonzero \(P_\alpha\in K[T_1,\ldots,T_s]\).  Fix
\(\alpha\) with \(P_\alpha\neq0\), and factor out a coefficient
\(c_\alpha\in K^\times\) of maximal absolute value:
\(P_\alpha=c_\alpha Q_\alpha\), where \(Q_\alpha\) has coefficients in
\(K^\circ\), at least one of absolute value \(1\).  The reduction
\(\widetilde Q_\alpha\in\widetilde K[T]\) is nonzero, and the residues
\(\tilde g_j\) are algebraically independent, so
\(\widetilde Q_\alpha(\tilde g)\neq0\); hence
\(w(Q_\alpha(g))=0\) and
\[
 w\bigl(f^\alpha P_\alpha(g)\bigr)
 =\langle\alpha,\gamma\rangle+v(c_\alpha).
\]
For \(\alpha\neq\alpha'\), equality of these values would give
\(\langle\alpha-\alpha',\gamma\rangle=v(c_{\alpha'})-v(c_\alpha)
\in\Gamma_K\), contradicting the independence of the \(\gamma_i\)
modulo \(\Gamma_K\otimes\mathbf Q\).  The finitely many nonzero
summands therefore have pairwise distinct values, and each is a nonzero
element of the field \(F\), of finite value.  A finite sum with a
unique term of strictly minimal value cannot vanish, by
\cref{lem:strong-triangle-equality}.  This contradiction proves the
independence.
\end{proof}

\begin{proposition}
\label{prop:berk-type-invariants}
\index{Berkovich affine line!type invariants}
For \(x\in\mathbf A^{1,\mathrm{an}}_K\), let
\(\widetilde{\mathcal H(x)}\) be the residue field and
\(|\mathcal H(x)^\times|\) the value group of the completed residue
field, and set
\[
 s(x)=\operatorname{trdeg}_{\widetilde K}\widetilde{\mathcal H(x)},
 \qquad
 t(x)=\dim_{\mathbf Q}
 \Bigl(\bigl(|\mathcal H(x)^\times|/|K^\times|\bigr)
 \otimes_{\mathbf Z}\mathbf Q\Bigr).
\]
Then \(s(x)+t(x)\leq1\), and:
\begin{enumerate}
\item \(x\) is of type I if and only if \(\mathcal H(x)=K\);
\item \(x\) is of type II if and only if \(s(x)=1\); in that case
      \(t(x)=0\), \(|\mathcal H(x)^\times|=|K^\times|\), and, for
      \(x=\zeta_{a,r}\) with \(r=|c|\),
      \(\widetilde{\mathcal H(x)}=\widetilde K(\tilde u)\) is a rational
      function field, generated by the residue \(\tilde u\) of
      \(u=(T-a)/c\);
\item \(x\) is of type III if and only if \(t(x)=1\); in that case
      \(s(x)=0\) and, for \(x=\zeta_{a,r}\),
      \(|\mathcal H(x)^\times|=|K^\times|\cdot r^{\mathbf Z}\);
\item \(x\) is of type IV if and only if \(s(x)=t(x)=0\) and
      \(\mathcal H(x)\neq K\); in that case \(\mathcal H(x)\) is a
      proper immediate extension of \(K\): it is complete with the same
      value group and the same residue field.
\end{enumerate}
In particular the type of a point is intrinsic.
\end{proposition}

\begin{proof}
Throughout, note that completion changes neither value group nor
residue field: if \(y_n\to y\neq0\) in a valued field, then eventually
\(|y_n-y|<|y|\), so \(|y_n|=|y|\) by the strong triangle equality; and
if \(|y|=1\), the elements \(y_n\) eventually have the same residue as
\(y\).  It therefore suffices to compute with the valued field
\(\bigl(K(T),|\cdot|_x\bigr)\) for points of support zero.  The
inequality \(s+t\leq1\) is \Cref{lem:berk-abhyankar-valued} with
\(F=K(T)\) and \(n=1\).

(1) If \(x\) is the evaluation at \(a\in K\), the character
\(K[T]\to K\), \(T\mapsto a\), identifies
\(\operatorname{Frac}(K[T]/(T-a))=K\), which is already complete, so
\(\mathcal H(x)=K\).  Conversely, if \(\mathcal H(x)=K\), the image
\(a\in K\) of \(T\) satisfies \(|f(x)|=|f(a)|\) for all \(f\), so
\(x\) is the evaluation at \(a\), of type I.

(2) Let \(x=\zeta_{a,r}\) with \(r=|c|\in|K^\times|\), and
\(u=(T-a)/c\).  Every value
\(|f(x)|=\max_i|c_i|r^i\) lies in \(|K^\times|\), and values of
quotients likewise, so \(|\mathcal H(x)^\times|=|K^\times|\) and
\(t(x)=0\).  For a polynomial \(P=\sum b_iS^i\in K^\circ[S]\),
\[
 |P(u)(x)|=\Bigl|\sum_ib_ic^{-i}(T-a)^i\Bigr|_x=\max_i|b_i|,
\]
so \(|P(u)(x)|=1\) exactly when \(\widetilde P\neq0\), and then the
residue of \(P(u)\) is \(\widetilde P(\tilde u)\), as one sees by
discarding the terms of absolute value less than one.  Hence
\(\tilde u\) is transcendental over \(\widetilde K\) and \(s(x)=1\).
For the reverse inclusion, let \(h=f/g\) with \(|h(x)|=1\).  Factor
\(f\) and \(g\) into linear factors.  A factor \(T-c_j\) with
\(|a-c_j|\leq r\) equals \(c\,(u-d_j)\) with
\(d_j=(c_j-a)/c\in K^\circ\), and \((u-d_j)\) has residue
\(\tilde u-\tilde d_j\in\widetilde K(\tilde u)\).  A factor \(T-c_j\)
with \(|a-c_j|>r\) equals
\((a-c_j)\bigl(1-(T-a)/(a-c_j)\bigr)\), whose second factor has
absolute value \(1\) at \(x\) and residue \(1\).  Multiplying the
normalizations, \(h\) is a constant of absolute value \(1\) times a
product of factors with residues in \(\widetilde K(\tilde u)\); its
residue lies in \(\widetilde K(\tilde u)\).  Thus
\(\widetilde{\mathcal H(x)}=\widetilde K(\tilde u)\).

(3) Let \(x=\zeta_{a,r}\) with \(r\notin|K^\times|\), \(r>0\).  Every
value \(\max_i|c_i|r^i\) lies in \(|K^\times|r^{\mathbf N}\), and the
maximum is attained at a unique index because
\(|c_i|r^i=|c_j|r^j\) with \(i\neq j\) would force
\(r^{i-j}\in|K^\times|\).  Values of quotients lie in
\(|K^\times|r^{\mathbf Z}\), and \(r=|(T-a)(x)|\) is a value, so
\(|\mathcal H(x)^\times|=|K^\times|r^{\mathbf Z}\) and
\(t(x)\geq1\).  By the Abhyankar inequality, \(t(x)=1\) and
\(s(x)=0\).

(4) Let \(x\) be of type IV with defining family \(E(x)\).  By Step 4
of the proof of \Cref{thm:berk-a1-classification}, for every
\(f\in K[T]\) there are \(a'\in K\) and a linear factorization with
\(|f(x)|=|u|\prod_j|c_j-a'|\in|K^\times|\); values of quotients
likewise, so \(t(x)=0\) and \(|\mathcal H(x)^\times|=|K^\times|\).
Next let \(f\in K[T]\) with \(|f(x)|=1\); we claim its residue lies in
\(\widetilde K\).  Choose \(a\in K\) beyond which the value of \(f\)
along \(E(x)\) has stabilized, and choose a strictly smaller disc
\(D_{a'}(x)\subsetneq D_a(x)\) with \(r_{a'}(x)<r_a(x)\), which exists
by Step 4.  Write \(f=\sum_ib_i'(T-a')^i\).  Since
\(\zeta_{a,r_a(x)}=\zeta_{a',r_a(x)}\) and
\(|f(\zeta_{a,r_a(x)})|=1\), we get \(|b_i'|\,r_a(x)^i\leq1\), so
\[
 |f-b_0'|_x=\Bigl|\sum_{i\geq1}b_i'(T-a')^i\Bigr|_x
 \leq\max_{i\geq1}\,r_a(x)^{-i}\,r_{a'}(x)^i
 =\frac{r_{a'}(x)}{r_a(x)}<1 .
\]
By the strong triangle equality, \(|b_0'|=1\) and \(f\) has the same
residue as the constant \(b_0'=f(a')\).  For a general element
\(h=f/g\) of absolute value \(1\), normalize \(f\) and \(g\) by
constants so that \(|f(x)|=|g(x)|=1\), which is possible because the
value group is \(|K^\times|\); then the residue of \(h\) is a quotient
of residues of constants, hence lies in \(\widetilde K\).  So
\(s(x)=0\), and \(\mathcal H(x)\) is an immediate extension of \(K\).
It is proper: \(\mathcal H(x)=K\) would make \(x\) of type I by (1),
contradicting exclusivity in \Cref{thm:berk-a1-classification}.

Finally, the four patterns
\(\bigl(\mathcal H(x)=K\bigr)\), \(\bigl(s=1\bigr)\),
\(\bigl(t=1\bigr)\), and
\(\bigl(s=t=0,\ \mathcal H(x)\neq K\bigr)\) are mutually exclusive and
cover the four types, so the type can be read off from
\(\mathcal H(x)\).
\end{proof}

\begin{definition}[Spherical completeness]
\index{spherically complete field}
A non-Archimedean field \(K\) is \emph{spherically complete} if every
chain of closed discs in \(K\), totally ordered by inclusion, has
nonempty intersection.  A chain of discs contains a countable
subfamily that is cofinal for reverse inclusion, because radii are real
numbers; so it is equivalent to impose the condition on decreasing
sequences of discs.
\end{definition}

\begin{proposition}
\label{prop:berk-type4-spherical}
Let \(K\) be complete, algebraically closed, and nontrivially valued.
Then \(\mathbf A^{1,\mathrm{an}}_K\) contains a point of type IV if and
only if \(K\) is not spherically complete.
\end{proposition}

\begin{proof}
If \(x\) is of type IV, then \(E(x)\) is a chain of closed discs with
empty intersection, so \(K\) is not spherically complete.

Conversely, suppose some chain of closed discs has empty intersection;
by the countable cofinality just observed we may take a decreasing
sequence \(D_n=D(a_n,r_n)\) with
\(\bigcap_nD_n=\varnothing\).  Define
\[
 |f(x)|:=\lim_{n\to\infty}|f(\zeta_{D_n})|
        =\inf_n|f(\zeta_{D_n})|,
\]
the limit of a nonincreasing bounded sequence
(\Cref{lem:berk-disc-combinatorics}).  Each defining property of a
point of \(\mathbf A^{1,\mathrm{an}}_K\) passes to the limit:
multiplicativity because a product of convergent sequences converges to
the product of the limits, the ultrametric inequality and the values on
\(K\) because weak inequalities and equalities survive limits.  So
\(x\in\mathbf A^{1,\mathrm{an}}_K\).

It remains to check \(\bigcap E(x)=\varnothing\); the classification
then makes \(x\) of type IV.  Suppose \(b\in\bigcap E(x)\).  For each
\(n\), and all \(m\geq n\), the center \(a_m\) lies in \(D_n\), so
\[
 r_{a_n}(x)=\lim_m\max\{|a_n-a_m|,r_m\}\leq r_n .
\]
Then \(|b-a_n|\leq r_{a_n}(x)\leq r_n\) shows \(b\in D_n\) for every
\(n\), contradicting \(\bigcap_nD_n=\varnothing\).
\end{proof}

\begin{example}[The four types over \(\mathbf C_p\)]
\label{ex:berk-types-cp}
Let \(\mathbf C_p\) be the completed algebraic closure of
\(\mathbf Q_p\); its value group is \(p^{\mathbf Q}\) and its residue
field is \(\overline{\mathbf F}_p\).  Then:
evaluation at any \(a\in\mathbf C_p\) is a type-I point;
\(\zeta_{0,1}\) is a type-II point, whose completed residue field has
residue field \(\overline{\mathbf F}_p(\tilde u)\) and value group
\(p^{\mathbf Q}\); and the point \(\zeta_{0,p^{-\lambda}}\) for
irrational \(\lambda\) is a type-III point, with value group
\(p^{\mathbf Q+\lambda\mathbf Z}\).  It is a classical fact, which we
neither prove nor use, that \(\mathbf C_p\) is not spherically
complete, so that type-IV points exist there as well.  The next
example carries out a completely explicit construction in equal
characteristic instead.
\end{example}

\begin{example}[An explicit type IV point]
\label{ex:berk-type4-puiseux}
Let \(K_0=\mathbf C(\!(t)\!)\) with \(v(t)=1\), let
\(P=\bigcup_{m\geq1}\mathbf C(\!(t^{1/m})\!)\) be the field of Puiseux
series, an algebraic closure of \(K_0\), and let \(K=\widehat P\) be its
completion, a complete algebraically closed field.

\emph{Description of \(K\).}  Every element of \(K\) is a series
\(y=\sum_{i\geq0}c_it^{q_i}\) with \(c_i\in\mathbf C^\times\) and
\(q_i\in\mathbf Q\) strictly increasing, the sequence being finite or
satisfying \(q_i\to+\infty\); and every such series converges in
\(K\).  Indeed, the partial sums of such a series lie in \(P\) and form
a Cauchy sequence, so the series converges.  Conversely, let
\(y\in K\) and choose \(y_n\in P\) with \(v(y-y_n)\geq n\).  For
\(q\in\mathbf Q\), the coefficient of \(t^q\) in \(y_n\) is independent
of \(n\) once \(n>q\): two Puiseux series whose difference has
valuation greater than \(q\) share all coefficients in degrees at most
\(q\).  Call this stable value \(c_q(y)\).  For fixed \(N\), the
support \(\{q\leq N:c_q(y)\neq0\}\) is contained in the support of a
single Puiseux series \(y_n\), \(n>N\), hence is finite.  Therefore
the support of \(y\) is a strictly increasing sequence, finite or
tending to \(+\infty\), and
\(v\bigl(y-\sum_{q\leq N}c_q(y)t^q\bigr)>N\) for every \(N\), so the
series converges to \(y\).

\emph{The nested discs.}  Define
\[
 x_n=\sum_{k=1}^{n}t^{\,1-\frac1k}
 =1+t^{1/2}+t^{2/3}+\cdots+t^{1-\frac1n},
 \qquad
 s_n=\rme^{-\left(1-\frac1{n+1}\right)},
\]
so \(s_n=|t^{1-\frac1{n+1}}|\) is strictly decreasing with limit
\(\rme^{-1}>0\).  Since \(|x_{n+1}-x_n|=|t^{1-\frac1{n+1}}|=s_n\), the
discs \(D_n=D(x_n,s_n)\) satisfy \(D_{n+1}\subseteq D_n\).

\emph{Empty intersection.}  Suppose \(y\in\bigcap_nD_n\).  Then
\(v(y-x_n)\geq1-\frac1{n+1}\) for every \(n\), so the series of \(y\)
agrees with \(x_n\) in all degrees \(q<1-\frac1{n+1}\); in particular
the coefficient of \(t^{1-\frac1k}\) in \(y\) equals \(1\) for every
\(k\leq n\).  As \(n\) is arbitrary, the support of \(y\) contains the
infinite set \(\{1-\frac1k:k\geq1\}\), which is bounded above by \(1\).
This contradicts the description of \(K\): an infinite support must
tend to \(+\infty\).  Hence \(\bigcap_nD_n=\varnothing\), \(K\) is not
spherically complete, and \Cref{prop:berk-type4-spherical} produces
from \((D_n)\) an explicit type-IV point of
\(\mathbf A^{1,\mathrm{an}}_K\), namely
\(|f(x)|=\lim_n|f(\zeta_{D_n})|\).
\end{example}

We now assemble the tree structure.  Inclusion of discs gives a partial
order on Gauss points, by \Cref{lem:berk-disc-combinatorics}; the
following proposition extends the picture to all points.

\begin{proposition}
\label{prop:berk-a1-tree}
\index{Berkovich affine line!tree structure}
Order \(\mathbf A^{1,\mathrm{an}}_K\) by \(x\leq y\) as above.
\begin{enumerate}
\item For every \(x\), the set \(A_x=\{y:y\geq x\}\) is totally
      ordered.  For a point \(x\) of type I, II, or III with disc
      \(D(b,r)\), one has \(A_x=\{\zeta_{b,s}:s\geq r\}\); for a point
      \(x\) of type IV, one has
      \(A_x=\{x\}\cup\{\zeta_{b,s}:b\in K,\ s\geq r_b(x)\}\), and
      \(A_x\setminus\{x\}\) consists of Gauss points of radii filling
      \((r_\infty(x),\infty)\).  In all cases the radius map is an
      order isomorphism from \(A_x\) onto \([\,r(x),\infty)\), where
      \(r(x)=\min_ar_a(x)\) in the first case and
      \(r(x)=r_\infty(x)\) in the second.
\item Points of types I and IV are minimal: \(y\leq x\) implies
      \(y=x\).
\item Any two points \(x,y\) have a least upper bound \(x\vee y\).
      If \(x=\zeta_{a,r}\) and \(y=\zeta_{b,u}\) are of types I--III,
      then
      \[
       x\vee y=\zeta_{a,m},\qquad m=\max\{r,u,|a-b|\}.
      \]
\end{enumerate}
\end{proposition}

\begin{proof}
We first record two facts.  For any \(b,u\),
\begin{equation}
 \zeta_{b,u}\geq x
 \iff u\geq r_b(x).
 \tag{4.1}\label{eq:berk-upper-bound-criterion}
\end{equation}
Indeed, if \(\zeta_{b,u}\geq x\) then
\(u=|(T-b)(\zeta_{b,u})|\geq r_b(x)\); conversely, if
\(u\geq r_b(x)\), then
\(x\leq\zeta_{b,r_b(x)}\leq\zeta_{b,u}\), by Step 1 of the proof of
\Cref{thm:berk-a1-classification} and by
\Cref{lem:berk-disc-combinatorics}.  Second, a strict upper bound of
any point is a Gauss point: suppose \(y\geq x\) and \(y\) is of type I
or IV.  If \(y\) is the evaluation at \(b\), then
\(r_b(x)\leq r_b(y)=0\), so \(x\) is the evaluation at \(b\) and
\(x=y\).  If \(y\) is of type IV, let \(c\in K\) be arbitrary; since
\(\bigcap E(y)=\varnothing\), there is \(a\) with
\(|c-a|>r_a(y)\geq r_a(x)\), and the strong triangle equality gives
\(r_c(x)=|c-a|=r_c(y)\).  Thus \(x\) and \(y\) agree on all linear
polynomials, hence on all of \(K[T]\) by multiplicativity and
factorization: \(x=y\).  This proves (2), and shows
\(A_x\subseteq\{x\}\cup\{\text{Gauss points}\}\).

(1) By \eqref{eq:berk-upper-bound-criterion},
\(A_x\setminus\{x\}\subseteq\{\zeta_{b,u}:u\geq r_b(x)\}\), with
equality of sets as stated; for \(x\) of type I--III with disc
\(D(b,r)\), the criterion \(u\geq r_{b'}(x)=\max\{|b-b'|,r\}\) says
exactly \(D(b',u)\supseteq D(b,r)\), so
\(A_x=\{\zeta_{b,s}:s\geq r\}\).  Total order: let
\(\zeta_{b,u},\zeta_{b',u'}\in A_x\), so \(D(b,u)\supseteq D_b(x)\)
and \(D(b',u')\supseteq D_{b'}(x)\).  The members \(D_b(x)\) and
\(D_{b'}(x)\) of \(E(x)\) are nested, say
\(D_b(x)\subseteq D_{b'}(x)\); then \(D(b,u)\) and \(D(b',u')\) both
contain \(D_b(x)\), hence intersect, hence are nested by
\Cref{lem:berk-disc-combinatorics}, and Gauss points of nested discs
are comparable.  The radius is strictly monotone on \(A_x\): nested
discs with equal radii coincide.  Its image is
\([r(x),\infty)\) for types I--III.  For type IV, every
\(u>r_\infty(x)\) occurs: choose \(a\) with \(r_a(x)<u\); then
\(\zeta_{a,u}\in A_x\).  The radius \(r_\infty(x)\) itself does not
occur: a disc \(D(b,r_\infty(x))\) containing some \(D_a(x)\) would
have radius at least \(r_a(x)>r_\infty(x)\).  Assigning to \(x\) the
parameter \(r(x)\) completes the order isomorphism in both cases.

(3) The set of common upper bounds is
\(A_x\cap A_y\), which is totally ordered by (1) and upward closed.
It is nonempty: by
\eqref{eq:berk-upper-bound-criterion}, \(\zeta_{b,u}\) with any
\(b\) and \(u\geq\max\{r_b(x),r_b(y)\}\) works.  If \(x=y\) the least
upper bound is \(x\).  Assume \(x\neq y\); we must produce a minimum.
For each \(b\in K\), the smallest upper bound with center \(b\) is
\(\zeta_{b,u_b}\) with \(u_b=\max\{r_b(x),r_b(y)\}\); set
\(\sigma=\inf_bu_b\), and choose \(b_n\) with
\(u_{b_n}\downarrow\sigma\).  The Gauss points
\(\zeta_{b_n,u_{b_n}}\) are pairwise comparable, so the discs
\(D(b_n,u_{b_n})\) contain a nested cofinal subsequence with radii
decreasing to \(\sigma\).  If some \(b_*\in\bigcap_nD(b_n,u_{b_n})\),
then
\[
 r_{b_*}(x)\leq\max\{r_{b_n}(x),|b_*-b_n|\}\leq u_{b_n}
 \longrightarrow\sigma,
\]
and similarly for \(y\), so \(u_{b_*}\leq\sigma\); by definition of
\(\sigma\), \(u_{b_*}=\sigma\), and
\(\zeta_{b_*,\sigma}\) is an upper bound of minimal radius, hence the
minimum of the totally ordered set \(A_x\cap A_y\).  If the
intersection is empty, the chain \((D(b_n,u_{b_n}))\) defines, by the
recipe of \Cref{prop:berk-type4-spherical}, a point \(z\) with
\(|f(z)|=\inf_n|f(\zeta_{b_n,u_{b_n}})|\geq|f(x)|\) for every \(f\),
i.e.\ a type-IV upper bound of \(x\) and of \(y\).  By (2) and the
second fact above, a type-IV upper bound forces \(z=x\) and \(z=y\),
contradicting \(x\neq y\).  So the minimum exists.

For the explicit formula, let \(x=\zeta_{a,r}\), \(y=\zeta_{b,u}\), and
\(m=\max\{r,u,|a-b|\}\).  Then
\(r_a(y)=\max\{|a-b|,u\}\le m\) and \(r_a(x)=r\leq m\), so
\(\zeta_{a,m}\) is an upper bound by
\eqref{eq:berk-upper-bound-criterion}.  Any upper bound
\(\zeta_{c,w}\) satisfies \(w\geq r_c(x)=\max\{|c-a|,r\}\) and
\(w\geq r_c(y)=\max\{|c-b|,u\}\); hence \(a\in D(c,w)\), \(w\geq r\),
\(w\geq u\), and \(w\geq|a-b|\) by the ultrametric inequality
\(|a-b|\le\max\{|c-a|,|c-b|\}\); so
\(D(a,m)\subseteq D(a,w)=D(c,w)\) and
\(\zeta_{a,m}\leq\zeta_{c,w}\).
\end{proof}

Two points have a least common
disc above them, and the chains between points supply
\(\mathbf P^{1,\mathrm{an}}_K\) with the structure of a topological tree.
The projective line is obtained by adding the classical point at infinity;
it is compact and uniquely path connected. This tree is the
one-dimensional prototype for the skeleta of \cref{chap:skeleta}.  A small
portion of the tree, including the passage from affine to projective space,
is shown in \cref{fig:berk-p1-tree}.

\begin{figure}[htbp]
\centering
\begin{tikzpicture}[x=1cm,y=1cm,>=Stealth,
  discpoint/.style={circle,draw,fill=white,inner sep=1.8pt},
  endpoint/.style={circle,draw,fill=black,inner sep=1.5pt}]
  \coordinate (join) at (0,1.05);
  \coordinate (za) at (-1.45,0.15);
  \coordinate (zb) at (1.45,0.15);
  \coordinate (a) at (-2.65,-0.95);
  \coordinate (b) at (2.65,-0.95);
  \coordinate (inf) at (0,3.05);
  \coordinate (ziv) at (3.15,1.75);
  \coordinate (xiv) at (4.15,2.35);

  \draw[very thick] (a)--(za)--(join)--(zb)--(b);
  \draw[very thick] (join)--(inf);
  \draw[very thick] (join)--(ziv);
  \draw[very thick,dashed] (ziv)--(xiv);

  \node[discpoint,label={[font=\small]above left:$\zeta_D$}] at (join) {};
  \node[discpoint,label={[font=\small]above left:$\zeta_{a,r}$}] at (za) {};
  \node[discpoint,label={[font=\small]above right:$\zeta_{b,s}$}] at (zb) {};
  \node[endpoint,label={[font=\small]below left:$a$}] at (a) {};
  \node[endpoint,label={[font=\small]below right:$b$}] at (b) {};
  \node[endpoint,label={[font=\small]above:$\infty$}] at (inf) {};
  \node[discpoint] at (ziv) {};
  \node[rectangle,draw,fill=white,inner sep=1.8pt,
    label={[font=\small]right:$x_{\mathrm{IV}}$}] at (xiv) {};

  \draw[-{Stealth[length=2mm]},thin]
    (-2.18,-0.65) to[bend left=12] (-0.35,0.83);
  \node[font=\scriptsize,fill=white,inner sep=1pt] at (-1.45,0.78)
    {increasing discs};
  \node[font=\scriptsize,align=left] at (3.55,1.12)
    {decreasing discs\\with empty intersection};
\end{tikzpicture}
\caption{A schematic finite subtree of \(\mathbf P^{1,\mathrm{an}}_K\).
The disc point \(\zeta_D\) is the least common disc above the two displayed
affine branches.  Solid circular endpoints represent type-I points; the square
endpoint of the dashed terminal branch represents a possible type-IV limit.
The drawing records incidence, not edge lengths.}
\label{fig:berk-p1-tree}
\end{figure}

\begin{example}[The standard annulus]\label{ex:berk-standard-annulus}
Assume \(K\) is discretely valued with \(v(\pi)=1\). In
\[
 \{x\in\mathbf A^{1,\mathrm{an}}_K:|\pi|\leq |T(x)|\leq1\},
\]
the seminorms
\[
 \left|\sum_i a_iT^i\right|_t
 =\max_i |a_i|\rme^{-it},\qquad 0\leq t\leq1,
\]
form a closed interval. The coordinate \(t=v_t(T)\) is affine along it.
The annulus retracts onto this interval; open discs are attached along its
points. The local model \(xy=\pi\) identifies \(t=v(x)\) and
\(1-t=v(y)\).
\end{example}

\section{Functoriality, base change, and GAGA}
\label{sec:berk-functoriality}

A morphism \(f:X\to Y\) of finite-type \(K\)-schemes induces
\[
 f^{\mathrm{an}}:X^{\mathrm{an}}\longrightarrow Y^{\mathrm{an}},
 \qquad
 |g(f^{\mathrm{an}}(x))|=|f^\#g(x)|.
\]
This construction respects identities and compositions. Closed and open
immersions, finite morphisms, and proper morphisms analytify to their
analytic counterparts.

Let \(L/K\) be a complete extension. Base change gives a continuous map
\[
 \pi_{L/K}:(X_L)^{\mathrm{an}}\longrightarrow X^{\mathrm{an}}.
\]
It is surjective. The analytic fiber over \(x\in X^{\mathrm{an}}\) is
canonically the spectrum of the completed scalar extension of the completed
residue field:
\[
              \pi_{L/K}^{-1}(x)\simeq
              \mathcal M\bigl(\mathcal H(x)\widehat\otimes_K L\bigr).
\]
The formula warns us that a base-changed point need not be unique. It also
explains why completed residue fields, rather than only algebraic residue
fields, belong to the basic language.

For proper schemes, Berkovich GAGA identifies coherent algebraic and
analytic sheaves and their cohomology. In particular, algebraic line
bundles and morphisms between proper schemes can be recovered from their
analytifications. We cite these results from \cite[Section~3.4]{Berk12}; the
metric constructions of \cref{chap:metrics} use them implicitly.

\section{The trivially valued case}
\label{sec:berk-trivial}

Let \(k\) now carry the trivial absolute value. Nothing in the affine
definition of \(X^{\mathrm{an}}\) changes. A point is equivalently a real
semivaluation
\[
 v_x:A\longrightarrow\mathbf R\cup\{+\infty\}
\]
that vanishes on \(k^\times\). The support
\(\{v_x=+\infty\}\) must be retained; it can be a nonzero prime ideal.

For \(t>0\), rescaling defines another point with the same support:
\[
                v_{t\cdot x}(f)=t\,v_x(f),
       \qquad |f(t\cdot x)|=|f(x)|^t.
\]
This is a continuous \(\mathbf R_{>0}\)-action. Along a fixed ray, the limit
as \(t\downarrow0\) is the trivial absolute value on
\(\kappa(\rho_X(x))\). More explicitly, as a seminorm on \(A\), the limit is
\(0\) on the fixed support \(\mathfrak p_x\) and \(1\) on
\(A\setminus\mathfrak p_x\); after passage to
\(\kappa(\mathfrak p_x)\), it is an absolute value and has zero kernel.

\begin{remark}[Cone terminology needs qualification]
If \(x\) has support zero on an integral affine scheme, its ray ends at the
trivial valuation of the function field. If the support is nonzero, it ends
at the trivial valuation of a proper subvariety. Thus the whole
analytification is stratified by cone-like pieces; it is not generally one
cone with one vertex. Moreover \(0\cdot(+\infty)\) is never used as an
extended-real arithmetic operation. The endpoint is defined by its
multiplicative seminorm.
\end{remark}

\begin{example}[The trivially valued affine line]
For \(r>0\), the formula
\[
 \left|\sum_i a_iT^i\right|_r=\max_{a_i\neq0}r^i
\]
defines a point of \(\mathbf A^{1,\mathrm{an}}_k\). In additive notation,
\(v_r(T)=-\log r\). The points with \(r<1\) have center at the origin on
\(\mathbf A^1_k\). For \(r=1\), the seminorm is trivial on \(k(T)\) and its
center is the generic point of \(\mathbf A^1_k\). When \(r>1\), the point
has no center on the affine model and is centered at infinity after passing
to \(\mathbf P^1_k\). This elementary example already shows why properness
matters for centers.
\end{example}

If \(X\) is proper over \(k\), \Cref{thm:berk-proper-compact} still
makes \(X^{\mathrm{an}}\) compact. What disappears is not the analytic
space, but the distinguished scale \(v(\pi)=1\). Homogeneity replaces that
normalization in the Boucksom--Jonsson theory of \cref{part:bj}.

\section{Dictionary established in this chapter}

The following correspondences will be used without further comment:
\[
\begin{array}{ccl}
\text{analytic point }x
 &\longleftrightarrow&
 \text{seminorm }|\cdot|_x
 \longleftrightarrow \text{ semivaluation }v_x,\\[1mm]
\text{support of }x
 &\longleftrightarrow&
 \ker|\cdot|_x\in\operatorname{Spec}A,\\[1mm]
\text{component }E\subset\mathcal X_s
 &\longleftrightarrow&
 x_E\text{ with }v_{x_E}=b_E^{-1}\operatorname{ord}_E,\\[1mm]
\text{proper algebraic variety}
 &\longrightarrow&
 \text{compact Hausdorff analytic space}.
\end{array}
\]
The next chapter places many divisorial and monomial points into finite
polyhedral subsets and constructs canonical retractions onto them.

\section*{Source notes}

The spectrum, analytic gluing, properness theorem, completed residue fields,
and GAGA are due to Berkovich; our principal reference is \cite{Berk12}.
Reduction maps and divisorial points in the precise model-theoretic form used
later are reviewed in \cite[Sections~2--3]{BFJ16}. The description of
\(\mathbf P^{1,\mathrm{an}}\) is standard; a supplementary source is
Baker--Rumely's monograph. This chapter deliberately separates the elementary
point-set construction from the deeper sheaf-theoretic results that are
cited.
